% concatenated from autosam.tex
% ============================================================================
% Automatica (autart) - Journal paper skeleton (OPTIMIZED PREAMBLE)
% - Minimal, template-friendly packages
% - No layout packages (e.g., geometry), no custom formatting
% - Theorem environments via amsthm (commonly accepted)
% - Structure + bullet-point placeholders
% ============================================================================

\documentclass[twocolumn]{autart}

% --- Minimal, standard packages (Automatica-friendly) ---
\usepackage{graphicx}    % figures
\usepackage{subcaption}
\usepackage{amsmath}     % math
\usepackage{amssymb}     % symbols
%\usepackage{amsthm}      % theorem environments
%\usepackage{amsmath,amssymb}
\usepackage{url}
\usepackage{enumitem}
\usepackage{xcolor}


\usepackage{comment}
\usepackage{algorithm}
\usepackage{algpseudocode}
\allowdisplaybreaks[2]


% Continuation line aligned with the current indentation level (algpseudocode)
\algnewcommand{\StateCont}[1]{\Statex\hspace{\algorithmicindent}#1}

\begin{document}

\begin{frontmatter}

\title{Decentralized Contingency MPC based on Safe Sets for Nonlinear Multi-agent Collision Avoidance}

% --- Replace with your authors/affiliations (keep it simple) ---
\author[Inst1]{Max Studt$^\ast$}\ead{m.studt@uni-luebeck.de}, \
\author[Inst1]{Georg Schildbach}\ead{georg.schildbach@uni-luebeck.de}
\address[Inst1]{Institute for Electrical Engineering in Medicine, University of Luebeck, Luebeck, Germany}

\begin{keyword}
Decentralized MPC; Contingency MPC; Nonlinear systems; Multi-agent systems; Collision avoidance;
Recursive feasibility; Safe sets; Lyapunov constraint
\end{keyword}

\begin{abstract}
Decentralized collision avoidance remains challenging, particularly when agents do not communicate any information related to planned trajectories. Most existing approaches either rely on conservative coordination mechanisms or provide limited guarantees on recursive feasibility and convergence. This paper develops a decentralized contingency MPC framework for multi-agent systems with nonlinear dynamics that achieves collision-free motion under a state-only information pattern. Each agent follows the same consensual rule set, enabling safe decentralized planning without communication. Each agent solves a local optimization problem that couples a nominal trajectory with a contingency certificate ensuring a feasible backup maneuver under receding-horizon operation. A novel geometric and decentralized safe-set update mechanism prevents feasibility loss between consecutive time steps. The resulting scheme guarantees recursive feasibility, including collision avoidance, and establishes a Lyapunov-type convergence result to an admissible safe equilibrium. Simulation results demonstrate performance in both sparse and dense multi-agent environments, including cluttered bottleneck scenarios and under plug-and-play operation.
\end{abstract}

\end{frontmatter}

% ============================================================================
% Theorem environments (define AFTER frontmatter; keep minimal)
% ============================================================================
\theoremstyle{plain}
\newtheorem{theorem}{Theorem}
\newtheorem{lemma}{Lemma}
\newtheorem{proposition}{Proposition}

\theoremstyle{definition}
\newtheorem{assumption}{Assumption}
\newtheorem{remark}{Remark}
\newtheorem{definition}{Definition}
%\newtheorem{proof}{Proof}
\newtheorem{corollary}{Corollary}
%\newenvironment{proof}{\par\noindent\textbf{Proof: }\ignorespaces}{\hfill$\square$\par}
\newenvironment{proof}
  {\par\noindent\textbf{\textit{Proof: }}\itshape}
  {\hfill$\square$\par}
% ============================================================================
\section{Introduction}
% ============================================================================

Multi-agent navigation and control problems require \emph{hard} collision avoidance while simultaneously achieving mission objectives such as reference tracking.
Model Predictive Control (MPC) provides a natural framework for this setting, since constraints can be handled explicitly while optimizing performance over a receding horizon~\cite{Mayne2000}.
\subsection{Literature overview}
\subsubsection*{Centralized planning}
A first class of approaches enforces collision avoidance in a \emph{centralized} manner by solving a coupled optimization problem for the entire fleet. Such formulations can encode global constraints directly, but typically scale poorly and require a coordinator with access to global information. Mixed-integer linear programs (MILPs) are a prominent example, enabling explicit collision-avoidance logic at the price of combinatorial complexity~\cite{Schouwenaars2001ECC,RichardsSchouwenaarsHowFeron2002}. Receding-horizon MILP/MPC variants have also been demonstrated in aerospace guidance applications, highlighting both the potential and computational burden of centralized approaches~\cite{SchouwenaarsValentiFeronHow2005}.

\subsubsection*{Distributed / decentralized MPC with coordination}
To improve scalability, many distributed or decentralized MPC schemes let agents solve local problems while coordinating to handle coupling constraints~\cite{Scattolini2009,Negenborn2009}. This coordination may rely on repeated communication rounds~\cite{CommunicationRounds}, negotiation of predicted trajectories~\cite{Negotiation}, consensus/agreement protocols~\cite{Lunze2022}, or distributed optimization methods such as ADMM/ALADIN~\cite{Aladin}. Related non-iterative approaches exchange compact descriptions of neighbor behavior, e.g., contracts of guaranteed future coupling-variable trajectories, which can support recursive feasibility and stability while limiting communication to local neighbors~\cite{LUCIA2015205}. However, such schemes still rely on communication and may introduce latency or coordination assumptions that are undesirable in safety-critical settings~\cite{Stewart2010,KOHLER20191,Venkat2007}.

An intermediate regime considers limited or time-varying communication, including bounded ranges, changing neighbor sets, and plug-and-play operation. Recent MPC formulations address this by reasoning about topology changes over the prediction horizon, for example via multi-trajectory MPC schemes that couple a nominal trajectory with a worst-case topology-change trajectory~\cite{Saccani2023}. More broadly, plug-and-play concepts have been studied in networked control and MPC for systems with changing components~\cite{Stoustrup2009,Riverso2013,LUCIA2015205}.

Several MPC-based approaches further reduce coordination effort toward single exchanges or local sensing. Examples include decentralized cooperative and reactive MPC-style navigation in unknown environments~\cite{HoyMatveevSavkin2012}, multi-robot NMPC with feasibility and stability analyses~\cite{LafmejaniBerman2021}, and decentralized deadlock-prevention layers that trigger coordination maneuvers in dense settings~\cite{MerryGoRound}. 

\subsubsection*{Decentralized approaches without communication}
In purely decentralized settings, agents observe only the current multi-agent configuration and cannot exchange planned trajectories or intentions. Geometric and reactive methods such as velocity obstacles~\cite{FioriniShiller1998} and ORCA~\cite{vanDenBergORCA2011} are computationally attractive and can provide collision avoidance under specific modeling assumptions. However, they are typically not designed to incorporate multi-step performance objectives, hard state/input constraints, and MPC-style closed-loop guarantees.

\subsubsection*{Safety filters and barrier-function based safety}
Optimization-based safety filters and control barrier functions enforce safety by rendering a safe set forward invariant via online correction problems~\cite{Ames2019}. Predictive safety filters connect this idea to MPC by embedding constraint satisfaction into a predictive structure~\cite{Wabersich2021}. However, combining multi-step planning, collision avoidance, and state-only decentralized information patterns with strong guarantees such as recursive feasibility and convergence remains challenging.

\subsection{Contributions}
To bridge these gaps, this paper adopts a \emph{contingency MPC} perspective. Contingency MPC couples a nominal plan with a safety or backup plan, thereby certifying the existence of an emergency maneuver under receding-horizon implementation~\cite{Alsterda2019}. A related communicationless multi-agent collision-avoidance setting has been addressed for \emph{linear} systems using pairwise bisecting-plane separation constraints and a deterministic backup evolution~\cite{Georg}. In contrast, this paper develops a decentralized contingency MPC framework for \emph{nonlinear} multi-agent systems under a state-only information pattern. The proposed method combines a local dual-plan formulation with a novel \textit{freeze-or-shift} (\textit{FoS}) safe-set update that preserves disjointness of active local safe regions over time. The resulting scheme provides tractable guarantees on recursive feasibility, collision avoidance, and Lyapunov-type convergence.

The main contributions of this paper are:
\begin{itemize}
  \item[C1] A fully decentralized \textit{FoS} update rule for time-varying convex local safe regions that preserves disjointness of the active safe sets.
  \item[C2] A decentralized contingency MPC formulation for nonlinear multi-agent systems that couples a nominal performance plan with a backup safety certificate via a shared first input.
  \item[C3] A constructive recursive-feasibility and collision-avoidance proof
based on shifted contingency candidates and invariant disjoint safe sets.
  \item[C4] A Lyapunov-type convergence result obtained by enforcing a monotone decrease condition on a contingency-related cost bound.
  \item[C5] A simulation study demonstrating effectiveness in sparse, dense, cluttered bottleneck, and plug-and-play multi-agent scenarios.
\end{itemize}

% ============================================================================
\section{Preliminaries and problem statement}
\label{sec:prelim}
% ============================================================================

\subsection{Notation}
\label{subsec:notation}
The sets of real and integer numbers are denoted by $\mathbb{R}$ and $\mathbb{Z}$, respectively. Moreover, $\mathbb{R}_+$ and $\mathbb{Z}_+$ denote the sets of nonnegative real and nonnegative integer numbers, respectively, and $\mathbb{Z}_{>0}$ denotes the set of positive integers. For integers $a,b\in\mathbb{Z}$ with $a\le b$, the (inclusive) index set is defined as
\(
\mathbb{Z}_a^b := \{k\in\mathbb{Z}\mid a \le k \le b\}.
\)
The set of agents is
\(
\mathcal{I} := \{1,\dots,M\}.
\)
Discrete time is indexed by $t\in\mathbb{Z}_+$ with sampling time $T_s>0$, while continuous time is denoted by $\tau\in\mathbb{R}_+$.
The next discrete time instant is denoted by $t^+ := t+1$ and the previous one by $t^- := t-1$.
For a vector $x\in\mathbb{R}^n$, $\|x\|$ denotes the Euclidean norm.
Let $n_p \in \mathbb{Z}_{>0}$ denote the dimension of the position space. For $c\in\mathbb{R}^{n_p}$ and $R\ge 0$, the closed Euclidean ball centered at $c$ with radius $R$ is defined as
\(
\mathbb{B}(c,R) := \{p\in\mathbb{R}^{n_p} : \|p-c\| \le R\}.
\)
The MPC prediction index $(k|t)$ is used throughout: $x_{i,(k|t)}$ denotes the prediction of $x_i(t+k)$ for agent $i$ computed at time $t$.
An equilibrium of agent $i$ is denoted by $(\bar x_i,\bar u_i)$ and satisfies
\(
\bar x_i = f_i(\bar x_i,\bar u_i).
\)
When needed, the associated equilibrium position is denoted by $\bar p_i := C_i \bar x_i$.
For the contingency plan, terminal safe equilibria are denoted by $(\bar x_i^{\mathrm c}, \bar u_i^{\mathrm c})$.

\subsection{Agent Dynamics and Constraints}
\label{subsec:dynamics}
Let \(\mathcal W \subseteq \mathbb{R}^{n_p}\) denote the workspace. For each agent \(i \in \mathcal{I}\), define the
body-feasible position set
\[
\mathcal W_i := \{p_i(t) \in \mathbb{R}^{n_p} : \mathbb{B}(p_i(t),r_i) \subseteq \mathcal W\}.
\]
Thus, \(p_i(t) \in \mathcal W_i\) means that the full body of agent \(i\), modeled as a
closed Euclidean ball of radius \(r_i\) centered at \(p_i(t)\), is contained in the workspace.
If no workspace-boundary constraints are present, one may simply take
\(\mathcal W= \mathcal W_i =\mathbb{R}^{n_p}\).

Each agent \(i \in \mathcal{I}\) is modeled by a discrete-time nonlinear
system
\begin{equation}
x_i(t^+) = f_i(x_i(t),u_i(t)),
\label{eq:agent_dynamics}
\end{equation}
with state \(x_i(t) \in \mathbb{R}^{n_i}\) and input
\(u_i(t) \in \mathbb{R}^{m_i}\). The position of agent \(i\) is given by
\(
p_i(t) = C_i x_i(t) \in \mathbb{R}^{n_p},
\)
where \(C_i \in \mathbb{R}^{n_p \times n_i}\) extracts the position
component of the state.

The sets of \emph{admissible states} and \emph{admissible inputs} for agent \(i\) are denoted by
\(
\mathcal X_i \subset \mathbb{R}^{n_i}, \ \mathcal U_i \subset \mathbb{R}^{m_i}.
\)
The closed-loop system trajectory is required to satisfy, for all $t \in \mathbb{Z}_+$
\(
x_i(t) \in \mathcal X_i, \ u_i(t) \in \mathcal U_i.
\)
It is assumed that the state constraints are chosen such that $p_i(t) \in \mathcal W_i$ whenever $x_i(t) \in \mathcal X_i$.

Each agent \(i\) is assigned an individual global reference state
\(
x_i^{\mathrm{ref}} \in \mathbb{R}^{n_i},
\)
with corresponding reference position
\(
p_i^{\mathrm{ref}} = C_i x_i^{\mathrm{ref}} \in \mathcal W_i.
\)
Throughout the paper, \(x_i^{\mathrm{ref}}\) is assumed to denote an
\emph{admissible equilibrium state} for agent \(i\), i.e., there exists an
admissible input \(u_i^{\mathrm{ref}} \in \mathcal U_i\) such that
\[
x_i^{\mathrm{ref}}
=
f_i(x_i^{\mathrm{ref}},u_i^{\mathrm{ref}}).
\]
\subsection{Collision Model and Safety Objective}
\label{subsec:collision}
Let 
\[
\mathcal B_i(t) := \mathbb{B}(p_i(t),r_i).
\]
denote the closed Euclidean ball occupied by agent $i$. A collision between two distinct agents $i\neq j$ at time $t$ occurs if
\[
\mathcal B_i(t)\cap \mathcal B_j(t)\neq\emptyset,
\]
equivalently if $\|p_i(t)-p_j(t)\| < r_i+r_j$.
The safety objective is collision-free closed-loop execution,
\[
\mathcal B_i(t)\cap \mathcal B_j(t)=\emptyset, \qquad \forall i\neq j,\ \forall t\in\mathbb{Z}_+,
\]
while driving each agent towards its reference or, if necessary, towards a safe admissible equilibrium.

\begin{assumption}[Exact nominal dynamics]
\label{ass:exact_dynamics}
The closed-loop evolution of each agent is exactly described by~\eqref{eq:agent_dynamics} under the input and state constraints $u_i(t)\in\mathcal U_i$ and $x_i(t)\in\mathcal X_i$, for all $t\in\mathbb{Z}_+$. In particular, no model mismatch and no further uncertainties are considered.
\end{assumption}

\begin{assumption}[Position invariance]
\label{ass:position_invariance}
The dynamics~\eqref{eq:agent_dynamics} are translation invariant. More precisely, for every agent $i$, there exists an embedding $\phi_i:\mathbb{R}^{n_p}\to\mathbb{R}^{n_i}$ such that $C_i\phi_i(\delta)=\delta$ and
\[
f_i(x_i+\phi_i(\delta),u_i)=f_i(x_i,u_i)+\phi_i(\delta)
\]
for all admissible $(x_i,u_i)$ and all $\delta\in\mathbb{R}^{n_p}$.
\end{assumption}

\begin{assumption}[Information pattern]
\label{ass:information_pattern}
At each time $t\in\mathbb{Z}_+$, agent $i\in\mathcal I$ can measure its own state $x_i(t)$ and the current states $x_j(t)$ of all other agents $j\in\mathcal I\setminus\{i\}$. The reference positions (or objectives) of other agents and their future inputs or predicted trajectories are not available.
\end{assumption}

% ============================================================================
\section{Local safe sets}
\label{sec:local_safe_sets}
% ============================================================================

This section introduces the geometric object used to enforce decentralized safety, namely the active local safe set associated with each agent. The active safe sets are chosen such that, from the current state of each agent, there exists a feasible backup maneuver that remains inside the set, reaches an admissible safe equilibrium in finite steps, and can subsequently be maintained there indefinitely.

At each time $t\in\mathbb{Z}_+$, each agent $i\in\mathcal{I}$ is associated with an active and time-varying local safe set
\begin{equation}
S_i^\ast(t) := \mathbb{B}\big(c_i(t),R_i(t)\big)\subseteq\mathbb{R}^{n_p},
\label{eq:local_safe_set}
\end{equation}
where $c_i(t)\in\mathbb{R}^{n_p}$ is the center and $R_i(t)\ge 0$ is the radius.
The active pair \((c_i(t),R_i(t))\) is determined by a deterministic
update rule based on a state-dependent safe-set generator. More
precisely, for each agent \(i\in\mathcal I\), let
\[
\Gamma_i:\mathcal X_i \rightarrow
\left\{
\mathbb B(c,R)\subseteq\mathbb R^{n_p}
\,:\,
c\in\mathbb R^{n_p},\ R\ge 0
\right\}
\]
denote a deterministic map that assigns to every admissible state
\(x_i\in\mathcal X_i\) a generated local safe set. For a given state
\(x_i\), this generated safe set is written as
\[
\Gamma_i(x_i)
=
\mathbb B\big(c_i(x_i),R_i(x_i)\big).
\]
The active safe set \(S_i^\ast(t)\) used by the controller need not
coincide with the generated set \(\Gamma_i(x_i(t))\). Instead, the
generated set may either be accepted or rejected by the \textit{FoS}
update rule introduced in Section~\ref{sec:FreezeandshiftOperator}.
Thus, reconstructing active safe sets generally requires the current
state information together with one-step memory of the previously
active safe-set parameters.
Independent of this update rule, the footprint of the agent must satisfy
\begin{equation}
\mathcal{B}_i(t)\subseteq S_i^\ast(t),\qquad \forall t\in\mathbb{Z}_+.
\label{eq:footprint_containment}
\end{equation}

\begin{assumption}[Local reconstructability]
\label{ass:local_reconstructability}
The deterministic safe-set generation and update rules are commonly
known to all agents. During closed-loop operation, each agent
\(i\in\mathcal I\) can reconstruct the active safe set \(S_j^\ast(t)\)
of every other agent \(j\in\mathcal I\setminus\{i\}\) from the observed
state \(x_j(t)\) and one-step memory of the previously active safe-set
parameters \(S_j^\ast(t^-)\), without communication of planned
trajectories or control inputs.
\end{assumption}

\begin{assumption}[Contingency recoverability]
\label{ass:contingency_recoverability}
For every agent $i\in\mathcal{I}$ and every time $t\in\mathbb{Z}_+$, the active safe set $S_i^\ast(t)$ is chosen such that there exists a feasible contingency maneuver starting from $x_i(t)$ that remains inside $S_i^\ast(t)$, satisfies all state and input constraints, and reaches an admissible safe equilibrium in finite time.
\end{assumption}

\begin{definition}[Admissible safe equilibrium]
\label{def:terminal_safe_equilibrium}
For every agent $i\in\mathcal{I}$ and every time $t\in\mathbb{Z}_+$ the state-input pair $(\bar x_i,\bar u_i)$ satisfying
\[
\bar x_i = f_i(\bar x_i,\bar u_i),
\qquad
\bar p_i := C_i\bar x_i \in S_i^\ast(t),
\]
and
\[
\mathbb{B}\big(\bar p_i,r_i\big)\subseteq S_i^\ast(t)\cap\mathcal{W}_i,
\]
represents an admissible safe equilibrium.
Moreover, once this equilibrium is reached, the agent can be kept inside $S_i^\ast(t)$ for all future times by applying the constant equilibrium input $\bar u_i$.
\end{definition}

Together, Assumption~\ref{ass:contingency_recoverability} and Definition~\ref{def:terminal_safe_equilibrium} formalize the intended role of the active safe set: It must support both finite-time recovery and indefinite safe holding. In particular, the active safe set must be chosen such that it contains a reachable admissible safe equilibrium and the agent can remain there indefinitely once this equilibrium has been reached.

\begin{remark}[Non-equilibrium-based terminal set]
\label{rem:terminal_invariant_set_extension}
The present paper considers systems for which terminal safety can be represented by an admissible safe equilibrium contained in the active safe set. For systems that do not admit such a fixed equilibrium, e.g., certain aircraft models, the same conceptual role could instead be played by a compact positively invariant terminal safe set $X_i^{\mathrm f}(t)\subseteq S_i^\ast(t)$ that is reachable in finite time and can be rendered invariant by a backup policy. %Such systems are not considered in the present paper.
\end{remark}

\begin{remark}[Model-dependent realization]
\label{rem:model_dependent_realization}
The concrete realization of Assumption~\ref{ass:contingency_recoverability} depends on the agent model. For some systems, local safe sets can be obtained from closed-form stopping bounds; for more general nonlinear systems, they may be constructed by reachable-set over-approximations, invariant-set arguments, or other conservative backup designs.
\end{remark}
% ============================================================================
\section{Decentralized contingency MPC formulation}
\label{sec:dcmpc}
% ============================================================================

This section states the decentralized contingency MPC problem solved by each agent $i\in\mathcal{I}$ at time $t\in\mathbb{Z}_+$.
At each time step, a single optimization simultaneously computes
(i) a nominal trajectory for performance and reference tracking and
(ii) a contingency trajectory that guarantees the existence of a feasible
backup maneuver within the local safe set.
The two plans are coupled by enforcing a shared first control input.

\subsection{Decision variables and horizons}
\label{subsec:variables_horizons}
Two prediction horizons are introduced: a nominal horizon $N_{\mathrm{n}}\in\mathbb{Z}_+$
and a contingency horizon $N_{\mathrm{c}}\in\mathbb{Z}_+$.
The nominal horizon $N_{\mathrm{n}}$ is chosen to shape tracking performance.
The contingency horizon $N_{\mathrm{c}}$ is selected sufficiently large such that a feasible contingency maneuver to an admissible safe equilibrium can be represented within the active local safe set.
Its choice depends on the agent dynamics, input and state constraints, and the construction of the local safe sets.

At time $t$, agent $i$ optimizes the following state--input sequences:

\textbf{Nominal trajectory:}
\[
X_i^\mathrm n(t) = \{x^{\mathrm{n}}_{i,(k|t)}\}_{k=0}^{N_{\mathrm{n}}}, \quad
U_i^\mathrm n(t) = \{u^{\mathrm{n}}_{i,(k|t)}\}_{k=0}^{N_{\mathrm{n}}-1}.
\]
\textbf{Contingency trajectory:}
\[
X_i^\mathrm c(t) = \{x^{\mathrm{c}}_{i,(k|t)}\}_{k=0}^{N_{\mathrm{c}}}, \quad
U_i^\mathrm c(t) = \{u^{\mathrm{c}}_{i,(k|t)}\}_{k=0}^{N_{\mathrm{c}}-1}.
\]
Each state $x^{\bullet}_{i,(k|t)} \in \mathbb{R}^{n_i}$ is associated with a position
\(
p^{\bullet}_{i,(k|t)} = C_i x^{\bullet}_{i,(k|t)} \in \mathbb{R}^{n_p},
\)
where $\bullet \in \{\mathrm{n},\mathrm{c}\}$ distinguishes nominal and contingency trajectories. In addition, the optimization includes the \emph{terminal contingency equilibrium pair} as a decision variable,
\begin{equation}
\bar{x}^{\mathrm{c}}_i(t) \in \mathcal X_i, \quad 
\bar{u}^{\mathrm{c}}_i(t) \in\mathcal{U}_i,  
\end{equation}
satisfying
\(
\bar{x}^{\mathrm{c}}_i(t) = f_i\big(\bar{x}^{\mathrm{c}}_i(t), \bar{u}^{\mathrm{c}}_i(t)\big).
\)
The closed-loop system evolves according to \eqref{eq:agent_dynamics} with the applied input given by the shared first move,
\begin{equation}
u_i(t) = u^{\mathrm{n},*}_{i,(0|t)} = u^{\mathrm{c},*}_{i,(0|t)},
\label{eq:shared_input_applied}
\end{equation}
where $(\cdot)^\ast$ denotes an optimal solution of the MPC problem at time $t$.

% ----------------------------------------------------------------------------
\subsection{Objective function}
\label{subsec:dcmpc_objective}
% ----------------------------------------------------------------------------
The local objective is kept generic, as the theoretical guarantees derived later rely on
structural properties rather than on a specific (e.g., quadratic) choice.
A typical design penalizes deviations of the nominal trajectory from a desired reference and regularizes the selection of the contingency equilibrium.

Let $\ell_i^{\mathrm{n}}(\cdot,\cdot)$ denote a tracking-related stage cost and
$V_i^{\mathrm{n}}(\cdot)$ a terminal tracking cost. Moreover, let
$V_i^{\mathrm{c}}(\bar{x}^{\mathrm{c}}_i(t), x_i^{\mathrm{ref}})$ denote an offset cost that measures
the distance of the contingency equilibrium to the desired target state.
\begin{equation}
\begin{aligned}
J_i\big(X_i^{\mathrm{n}}(t), &U_i^{\mathrm{n}}(t), \bar{x}^{\mathrm{c}}_i(t), x_i^{\mathrm{ref}}\big)
:= \\
\sum_{k=0}^{N_{\mathrm{n}}-1}
&\ell_i^{\mathrm{n}}\!\big(x^{\mathrm{n}}_{i,(k|t)}, u^{\mathrm{n}}_{i,(k|t)}\big)
\;+\;
V_i^{\mathrm{n}}\!\big(x^{\mathrm{n}}_{i,(N_{\mathrm{n}}|t)}, x_i^{\mathrm{ref}}\big) \\
&\;+\;
\gamma\,V_i^{\mathrm{c}}\!\big(\bar{x}^{\mathrm{c}}_i(t), x_i^{\mathrm{ref}}\big).
\end{aligned}
\label{eq:generic_objective}
\end{equation}
The scalar $\gamma>0$ weights the preference for contingency equilibria that are close to the desired reference. Note that the cost term related to the selected terminal contingency equilibrium state is required in order to enforce a Lyapunov-like decrease
through an additional constraint (Section~\ref{subsec:dcmpc_convergence_constraint}).

\begin{definition}[Optimal contingency equilibrium]
\label{def:opt_reachable_equilibrium}
At each time step $t\in\mathbb{Z}_+$, an \emph{optimal reachable contingency equilibrium} minimizes the equilibrium-offset cost, i.e.,
\[
(\bar x_i^{\mathrm c,\ast}(t),\bar u_i^{\mathrm c,\ast}(t))
\in
\arg\min_{(\bar x_i^{\mathrm c},\bar u_i^{\mathrm c})\in\mathcal Z_i^{\mathrm c}(t)}
V_i^{\mathrm c}\!\big(\bar x_i^{\mathrm c},x_i^{\mathrm{ref}}\big),
\]
where $\mathcal Z_i^{\mathrm c}(t)$ denotes the set of admissible terminal contingency equilibrium pairs.
\end{definition}
\begin{remark}[Role of the equilibrium-offset weight]
The scalar $\gamma>0$ regularizes the selection of the contingency equilibrium. In related multi-trajectory MPC formulations, it can be shown that for sufficiently large equilibrium-offset weights, the equilibrium selected by the optimizer can be made arbitrarily close to the offset-minimizing admissible equilibrium ~\cite[Lemma~1]{Saccani2023}.
\end{remark}

% ----------------------------------------------------------------------------
\subsection{Constraints}
\label{subsec:dcmpc_constraints}
% ----------------------------------------------------------------------------
Both predicted trajectories start from the measured state,
\begin{equation}
x^{\mathrm{n}}_{i,(0|t)} = x^{\mathrm{c}}_{i,(0|t)} = x_i(t).
\label{eq:init_both}
\end{equation}
For the nominal trajectory, the dynamics \eqref{eq:agent_dynamics} are enforced for
$k=0,\ldots,N_{\mathrm{n}}-1$,
\[
x^{\mathrm{n}}_{i,(k+1|t)} = f_i\big(x^{\mathrm{n}}_{i,(k|t)}, u^{\mathrm{n}}_{i,(k|t)}\big),
\]
and analogously for the contingency trajectory for $k=0,\ldots,N_{\mathrm{c}}-1$,
\[
x^{\mathrm{c}}_{i,(k+1|t)} = f_i\big(x^{\mathrm{c}}_{i,(k|t)}, u^{\mathrm{c}}_{i,(k|t)}\big).
\]
Stagewise state and input constraints are imposed by
\[
x^{\mathrm{n}}_{i,(k+1|t)}\in\mathcal{X}_i,\ \ u^{\mathrm{n}}_{i,(k|t)}\in\mathcal{U}_i,
\qquad \forall k\in\mathbb{Z}_0^{N_{\mathrm{n}}-1},
\]
\[
x^{\mathrm{c}}_{i,(k+1|t)}\in\mathcal{X}_i,\ \ u^{\mathrm{c}}_{i,(k|t)}\in\mathcal{U}_i,
\qquad \forall k\in\mathbb{Z}_0^{N_{\mathrm{c}}-1}.
\]

To ensure that the predicted contingency trajectory represents a valid backup plan, the terminal equilibrium constraint
\begin{equation}
x^{\mathrm{c}}_{i,(N_{\mathrm{c}}|t)} = \bar{x}^{\mathrm{c}}_i(t),
\qquad
\bar{x}^{\mathrm{c}}_i(t) = f_i\big(\bar{x}^{\mathrm{c}}_i(t), \bar{u}^{\mathrm{c}}_i(t)\big),
\label{eq:terminal_steady_state}
\end{equation}
is imposed.
The terminal position is required to satisfy
\[
\bar{p}^{\mathrm{c}}_i(t) := C_i\bar{x}^{\mathrm{c}}_i(t)\in S_i^\ast(t).
\]
Moreover, the terminal equilibrium must be reachable from the current state $x_i(t)$ within $N_{\mathrm{c}}$ steps by a feasible contingency trajectory that remains inside the active local safe set, i.e.,
\(
\bar{x}_i^c(t) \in \mathrm{reach}_{N_c}(x_i(t)).
\)

\subsubsection{Local Safe-Set Constraints}

The entire contingency position plan is required to remain inside the current active local safe set \(S_i^\ast(t)\). To account for the agent size, the constraint is imposed with a radius margin \(r_i\), i.e.,
\begin{equation}
\|p^{\mathrm{c}}_{i,(k|t)} - c_i(t)\| \le R_i(t) - r_i,
\qquad \forall k\in\mathbb{Z}_0^{N_{\mathrm{c}}}.
\label{eq:contingency_plan_in_active_safe_set}
\end{equation}
Equivalently, \(\mathbb{B}\bigl(p^{\mathrm{c}}_{i,(k|t)},r_i\bigr)\subseteq S_i^\ast(t)\) for all
\(k\in\mathbb{Z}_0^{N_{\mathrm{c}}}\).

Furthermore, for every prediction step \(k\), the remainder of the
contingency trajectory (from step \(k\) onward) is required to remain
inside the local safe set generated by \(\Gamma_i\) from the predicted
contingency state \(x^{\mathrm c}_{i,(k|t)}\). Let this generated safe set
be given by
\[
\Gamma_i\bigl(x^{\mathrm c}_{i,(k|t)}\bigr)
=
\mathbb B\bigl(
c_i(x^{\mathrm c}_{i,(k|t)}),
R_i(x^{\mathrm c}_{i,(k|t)})
\bigr).
\]
Then the tail-containment constraint is
\begin{equation}
\begin{aligned}
\|p^{\mathrm c}_{i,(l|t)} - c_i(x^{\mathrm c}_{i,(k|t)})\|
\le
R_i(x^{\mathrm c}_{i,(k|t)})-r_i, \\
\qquad
\forall k\in\mathbb Z_0^{N_{\mathrm c}},
\ \forall l\in\mathbb Z_k^{N_{\mathrm c}}.
\label{eq:tail_containment_constraint}
\end{aligned}
\end{equation}
Constraint~\eqref{eq:tail_containment_constraint} enforces a receding-horizon recoverability property:
At each prediction stage \(k\), the planned tail
\[
\left\{x^{\mathrm{c}}_{i,(l|t)}\right\}_{l=k}^{N_{\mathrm{c}}}
\]
remains entirely inside the candidate local safe set generated from the predicted state \(x^{\mathrm{c}}_{i,(k|t)}\).
In particular, after applying the first input, the shifted tail
\[
\left\{x^{\mathrm{c}}_{i,(l+1|t)}\right\}_{l=0}^{N_{\mathrm{c}}-1}
\]
is contained in the candidate safe set generated from the new initial condition
\(
x_i(t^+) = x^{\mathrm{c}}_{i,(1|t)}.
\)
Hence, if the active safe set at the next time step is chosen as this generated candidate safe set, the shifted contingency tail remains admissible with respect to the newly induced active safe set.
This property is a key ingredient for establishing recursive feasibility under receding-horizon implementation, including the case where the active safe set is updated by the \textit{FoS} rule.

Figure~\ref{fig:TailContainmentInfeasibility} illustrates why constraint~\eqref{eq:tail_containment_constraint} is necessary.
Even if the local problem is feasible at time $t$ with active safe set $S_i^\ast(t)$, recursive feasibility may be lost after applying the first input if the remaining contingency tail is not guaranteed to lie inside the safe set induced at the successor state.
\begin{figure}[h]
    \centering
    \includegraphics[width=0.4\textwidth]{Figures/Feasibility_Loss.png}
    \caption{Loss of recursive feasibility without constraint \eqref{eq:tail_containment_constraint}. Although the local MPC problem is feasible at time $t$ with active safe set $S^\ast(t)$ (red), the naive update to $t^+$ yields a set $S^\ast(t^+)$ (dashed) that no longer contains the complete tail of the contingency plan (blue) from previous time step. This violates the safe-set containment constraints and renders the next-step problems infeasible.}
    \label{fig:TailContainmentInfeasibility}
\end{figure}

\subsubsection{Collision avoidance w.r.t.\ static obstacles}
\label{subsec:CollisionAvoidance}
Collision avoidance with respect to static obstacles is enforced explicitly along the
contingency trajectory. Let $\mathcal{O}\subset\mathbb{R}^{n_p}$ denote the obstacle set (union of closed obstacle
regions) and define a continuous obstacle-avoidance function
$h:\mathbb{R}^{n_p}\rightarrow\mathbb{R}$ such that
\[
h(p_i(t))\ge 0 \ \Longleftrightarrow\ \mathcal{B}_i\cap \mathcal{O} = \emptyset.
\]
Then the following hard constraints are imposed for all predicted contingency positions:
\begin{equation}
h(p^{\mathrm{c}}_{i,(k|t)}) \ge 0,
\qquad \forall\, i\in\mathcal{I},\ \forall\, k\in\mathbb{Z}_0^{N_{\mathrm{c}}}.
\label{eq:static_obstacle_constraints}
\end{equation}

% ----------------------------------------------------------------------------
\subsubsection{Convergence constraint}
\label{subsec:dcmpc_convergence_constraint}
% ----------------------------------------------------------------------------
Following the contingency MPC rationale, a decrease in a contingency-cost functional is enforced through an inequality constraint.

Let $\ell_i^{\mathrm{c}}(\cdot,\cdot)$ denote a continuous, nonnegative contingency stage cost, and define
\begin{equation}
\begin{aligned}
J_i^{\mathrm{c}}(t)
:=
\sum_{k=0}^{N_{\mathrm{c}}-1}
&\ell_i^{\mathrm{c}}\!\big(x^{\mathrm{c}}_{i,(k|t)}-\bar{x}^{\mathrm{c}}_i(t),\,u^{\mathrm{c}}_{i,(k|t)}-\bar{u}^{\mathrm{c}}_i(t)\big) \\
&\;+\;
V_i^{\mathrm{c}}\!\big(\bar{x}^{\mathrm{c}}_i(t), x_i^{\mathrm{ref}}\big).
\end{aligned}
\label{eq:c_cost_def}
\end{equation}
A bound $\hat{J}_i^{\mathrm{c}}(t)\in\mathbb{R}_+$ is maintained recursively and the following constraint is imposed:
\begin{equation}
J_i^{\mathrm{c}}(t) \le \hat{J}_i^{\mathrm{c}}(t).
\label{eq:c_cost_decrease_constraint}
\end{equation}
The update law for $\hat{J}_i^{\mathrm{c}}(t)$ is constructed from the shifted tail of the previously optimal contingency trajectory and yields the one-step bound
\begin{equation}
\hat J_i^{\mathrm c}(t^+)
:=
J_i^{\mathrm c,\ast}(t)
-
\ell_i^{\mathrm{c}}\!\big(x_i(t)-\bar{x}_i^{\mathrm{c},\ast}(t),\,u_i(t)-\bar{u}_i^{\mathrm{c},\ast}(t)\big),
\label{eq:lyap_bound_update_main}
\end{equation}
where $x_i(t) = x_{i,(0|t)}^{c,\ast}$ and $u_i(t)=u^{\mathrm c,\ast}_{i,(0|t)}$ is the applied shared first input. The feasibility of this shifted-tail bound and its role in the convergence proof are established in the Appendix.

\begin{remark}[Equilibrium]
\label{rem:tv_ss_case_split}
Since the optimizer may select a different terminal contingency equilibrium state $\bar{x}^c_i(t)$ at different time steps, two situations can occur.

\begin{enumerate}
\item[(i)] \emph{Reference becomes reachable within the active safe set.}
There exists a time instant $\bar{t}\in\mathbb{Z}_+$ such that the global reference is feasible as a
contingency terminal state thereafter, i.e.,
\[
\exists\,\bar{t}\in\mathbb{Z}_+:\quad
\bar{x}^{c}_i(t)=x_i^{\mathrm{ref}} \qquad \forall t\ge \bar{t},
\]
In this case the selected terminal equilibrium  is constant for all $t\ge\bar{t}$ and the Lyapunov-type analysis
yields convergence to the reference state.

\item[(ii)] \emph{Reference is not reachable due to persistent blocking.}
There exists a time instant $\bar{t}\in\mathbb{Z}_+$ such that the reference cannot be selected as a feasible terminal state, and the optimal reachable terminal state remains constant thereafter, i.e.,
\[
\exists\,\bar{t}\in\mathbb{Z}_+,\ \exists\,\bar{x}^{c,\star}_i:\quad
\bar{x}^{c}_i(t)=\bar{x}^{c,\star}_i \qquad \forall t\ge \bar{t}.
\]
Then the convergence constraint implies convergence to the fixed terminal state $\bar{x}^{c,\star}_i$.

\end{enumerate}

If $\bar{x}^c_i(t)$ changes infinitely often, the constraint \eqref{eq:c_cost_decrease_constraint}
still enforces a monotonic decrease of the contingency-cost bound generated by the shifted-tail update,
but a standard ``convergence to a fixed equilibrium'' statement is no longer immediate. In this case,
the constraint is interpreted as enforcing progress towards the \emph{currently selected} contingency
terminal state and providing a stabilizing regularization of the contingency plan.
\end{remark}

% ----------------------------------------------------------------------------
\subsection{Local finite-horizon optimal control problem}
\label{subsec:dcmpc_problem_statement}
% ----------------------------------------------------------------------------
For each agent $i\in\mathcal{I}$ and time $t\in\mathbb{Z}_+$, a local finite-horizon optimal control problem (FHOCP) is solved:

\refstepcounter{equation}\label{prob:local_fhocp}\addtocounter{equation}{-1}
\begin{align}
& \min_{X_i^{\mathrm{n,c}},\,U_i^{\mathrm{n,c}}\,\bar{x}_i^{\mathrm{c}},\,\bar{u}_i^{\mathrm{c}}}\
  J_i\big(X_i^{\mathrm{n}}(t), U_i^{\mathrm{n}}(t), \bar{x}^{\mathrm{c}}_i(t), x_i^{\mathrm{ref}}\big) \tag{15a}
  \label{prob:local_fhocp_obj}
\\
& \text{subject to:} \notag
\\
& x^{\mathrm{n}}_{i,(0|t)} = x^{\mathrm{c}}_{i,(0|t)} = x_i(t) \tag{15b}
  \label{prob:local_fhocp_init}
\\
& u^{\mathrm{n}}_{i,(0|t)} = u^{\mathrm{c}}_{i,(0|t)} \tag{15c}
  \label{prob:local_fhocp_shared}
\\
& x^{\mathrm{n}}_{i,(k+1|t)} = f_i \big(x^{\mathrm{n}}_{i,(k|t)}, u^{\mathrm{n}}_{i,(k|t)}\big),
  \forall k\in\mathbb{Z}_0^{N_{\mathrm{n}}-1} \tag{15d}
  \label{prob:local_fhocp_dyn_n}
\\
& x^{\mathrm{c}}_{i,(k+1|t)} = f_i \big(x^{\mathrm{c}}_{i,(k|t)}, u^{\mathrm{c}}_{i,(k|t)}\big),
  \forall k\in\mathbb{Z}_0^{N_{\mathrm{c}}-1} \tag{15e}
  \label{prob:local_fhocp_dyn_c}
\\
& x^{\mathrm{n}}_{i,(k+1|t)}\in\mathcal{X}_i,\ \ u^{\mathrm{n}}_{i,(k|t)}\in\mathcal{U}_i,
  \forall k\in\mathbb{Z}_0^{N_{\mathrm{n}}-1} \tag{15f}
  \label{prob:local_fhocp_bounds_n}
\\
& x^{\mathrm{c}}_{i,(k+1|t)}\in\mathcal{X}_i,\ \ u^{\mathrm{c}}_{i,(k|t)}\in\mathcal{U}_i,
  \forall k\in\mathbb{Z}_0^{N_{\mathrm{c}}-1} \tag{15g}
  \label{prob:local_fhocp_bounds_c}
\\
& h \bigl(p^{\mathrm{c}}_{i,(k|t)}\bigr) \ge 0,
  \forall k\in\mathbb{Z}_0^{N_{\mathrm{c}}} \tag{15h}
  \label{prob:static_obstacle_constraintsMPC}
\\
& x^{\mathrm{c}}_{i,(N_{\mathrm{c}}|t)} = \bar{x}^{\mathrm{c}}_i(t) \tag{15i}
  \label{prob:local_fhocp_terminal_state}
\\
%& \bar{x}^{\mathrm{c}}_i(t) = f_i \bigl(\bar{x}^{\mathrm{c}}_i(t),\bar{u}^{\mathrm{c}}_i(t)\bigr) \tag{15j}
%  \label{prob:local_fhocp_terminal_equilibrium}
& \bar{x}^{\mathrm{c}}_i(t)
= f_i \bigl(\bar{x}^{\mathrm{c}}_i(t),\bar{u}^{\mathrm{c}}_i(t)\bigr),
\ \bigl(\bar{x}^{\mathrm{c}}_i(t), 
\bar{u}^{\mathrm{c}}_i(t)\bigr)\in \mathcal{X}_i \times \mathcal{U}_i
\tag{15j}
\label{prob:local_fhocp_terminal_equilibrium}
\\
& C_i\bar{x}^{\mathrm{c}}_i(t)\in S_i^\ast(t) \tag{15k}
  \label{prob:local_fhocp_terminal_position}
\\
& \|p^{\mathrm{c}}_{i,(k|t)} - c_i(t)\| \le R_i(t) - r_i,
  \forall k\in\mathbb{Z}_0^{N_{\mathrm{c}}} \tag{15l}
  \label{prob:local_fhocp_containment}
\\
& \|p^{\mathrm{c}}_{i,(l|t)} - c_i(x^{\mathrm{c}}_{i,(k|t)})\| \leq R_i(x^{\mathrm{c}}_{i,(k|t)}) - r_i, \tag{15m}
  \label{prob:local_fhocp_tail}
\\
& \qquad \qquad \qquad \forall k\in\mathbb{Z}_0^{N_{\mathrm{c}}},\ \forall l\in\mathbb{Z}_k^{N_{\mathrm{c}}} \notag
\\
& J_i^{\mathrm{c}}(t) \le \hat{J}_i^{\mathrm{c}}(t) \tag{15n}
  \label{prob:local_fhocp_lyap}
\end{align}
\setcounter{equation}{15}

Problem~\eqref{prob:local_fhocp} is, generally, a non-convex Nonlinear Program (NLP), even if the dynamics are linear. It is solved in receding-horizon fashion. The applied input is the shared first input
\eqref{prob:local_fhocp_shared}.

\begin{remark}[Performance and safety]
No explicit inter-agent collision-avoidance constraints are imposed on the nominal prediction
\(X_i^\mathrm n\). Safety is ensured by the contingency trajectory
\(X_i^\mathrm c\) via
\eqref{prob:local_fhocp_containment}--\eqref{prob:local_fhocp_tail}. The nominal part can therefore be used flexibly for performance shaping, e.g., through additional objectives or soft constraints; the theoretical guarantees rely only on the contingency-plan feasibility structure.
\end{remark}

%============================================================================
\section{\textit{FoS} update of the local safe sets}
\label{sec:FreezeandshiftOperator}
% ============================================================================

Recall that each agent $i$ is associated with an active local safe set as described in~\eqref{eq:local_safe_set}.
After solving the local FHOCPs at time \(t\) and applying the
shared first input, the closed-loop dynamics of each agent follow
\eqref{eq:agent_dynamics}.
The successor state of agent \(i\) is denoted by
\(
x_i(t^+) \in \mathcal X_i.
\)
Recall the deterministic state-dependent safe-set generator
\(\Gamma_i\) introduced in Section~\ref{sec:local_safe_sets}.
Given the successor state \(x_i(t^+)\), the candidate safe set for the
next time step is defined as
\begin{equation}
\tilde S_i(t^+)
:=
\Gamma_i\bigl(x_i(t^+)\bigr)
=
\mathbb{B}\bigl(\tilde c_i(t^+), \tilde R_i(t^+)\bigr),
\label{eq:candidate_safe_set}
\end{equation}
where \(\tilde c_i(t^+):=c_i(x_i(t^+))\) and \(\tilde R_i(t^+):=R_i(x_i(t^+))\). At time \(t\), agent \(i\) checks whether updating its active safe set
to the candidate \(\tilde S_i(t^+)\) would create an overlap with
either (i) the candidate safe set \(\tilde S_j(t^+)\) of another agent
\(j\neq i\), or (ii) the currently active safe set \(S_j^\ast(t)\) of
another agent \(j\neq i\). Accordingly, the \emph{freeze indicator} ($\chi$) for agent $i$ at time $t$ is defined as
\begin{equation}
\label{eq:freezeIndicator}
\chi_i(t)
:=
\begin{cases}
1 \ \text{if} & \exists j\in\mathcal{I}\setminus\{i\}: 
\begin{aligned}[t]
&\tilde{S}_i(t^+)\cap \tilde{S}_j(t^+)\neq\emptyset \\
\ \ &\text{or}\ \tilde{S}_i(t^+)\cap S_j^\ast(t)\neq\emptyset,
\end{aligned}
\\[0.2em]
0 & \text{otherwise}.
\end{cases}
\end{equation}
Since all local safe sets are Euclidean balls, the above condition is equivalent to the distance tests
\begin{align}
    \exists j \neq i:\ 
    &\|\tilde{c}_i(t^+) - \tilde{c}_j(t^+)\|
      < \tilde{R}_i(t^+) + \tilde{R}_j(t^+)
      \nonumber\\
      &\quad \text{or}\ 
      \|\tilde{c}_i(t^+) - c_j(t)\|
      < \tilde{R}_i(t^+) + R_j(t).
\label{eq:freeze_distance_test}
\end{align}
The \emph{FoS} update rule for the active local safe sets is then
defined as
\begin{equation}
\label{eq:freeze_shift_update}
(c_i(t^+), R_i(t^+))
:=
\begin{cases}
\big(c_i(t), R_i(t)\big), &\text{if } \chi_i(t)=1,\\[0.2em]
\big(\tilde{c}_i(t^+), \tilde{R}_i(t^+)\big), &\text{if } \chi_i(t)=0.
\end{cases}
\end{equation}
Equivalently,
\[
S_i^\ast(t^+)
=
\begin{cases}
S_i^\ast(t), &\text{if } \chi_i(t)=1,\\[0.2em]
\tilde S_i(t^+), &\text{if } \chi_i(t)=0.
\end{cases}
\]
\begin{remark}[Decentralized update]
\label{rem:decentralized_update}
By Assumption \ref{ass:information_pattern} and the deterministic construction of the candidate safe sets, each agent can locally reconstruct $S_j^\ast(t)$ and $\tilde S_j(t^+)$ for all $j\neq i$ from the current state and one-step memory of the previously active safe sets.
\end{remark}

\begin{figure}[h]
    \centering
    \includegraphics[width=0.45\textwidth]{Figures/Feasibility_LossNoVoronoi.png}
    \caption{Loss of recursive feasibility without \textit{FoS}. Although the local MPC problems are feasible at time $t$ with disjoint active safe sets $S_i^\ast(t)$ (red), the naive update to $t^+$ yields sets $S_i^\ast(t^+)$ (dashed) that are no longer pairwise disjoint. The resulting overlaps can violate the safe-set constraints and render the next-step problems infeasible.}
    \label{fig:FreezeShiftRule3Agents}
\end{figure}
Without the \textit{FoS} update rule, pairwise disjointness of the
active safe sets is not forward invariant. Even if
\(\{S_j^\ast(t)\}_{j\in\mathcal I}\) is pairwise disjoint, the naive
update \(S_i^\ast(t^+)=\Gamma_i(x_i(t^+))\), applied independently by
all agents, may yield overlapping safe sets for some \(i\neq j\); see
Fig.~\ref{fig:FreezeShiftRule3Agents}. Such overlaps can invalidate
the geometric separation used in the contingency constraints and may
lead to loss of recursive feasibility.
The \textit{FoS} rule prevents this by decoupling the model-dependent
generation of candidate safe sets from the geometric decision whether
to \textit{shift} the active safe set to the new candidate or to
\textit{freeze} it at its previously active value.


In the following, the closed-loop execution of the proposed decentralized contingency MPC scheme is given in algorithmic form.
\begin{algorithm}[H]
\caption{Decentralized Contingency MPC with \textit{FoS} Safe Sets}
\label{alg:dcmpc_freeze_shift}
\begin{algorithmic}[1]
\Require Agents \(M\), sampling time \(T_s\), horizons \(N_{\mathrm n},N_{\mathrm c}\), local constraints, $x_i^{\mathrm{ref}}$, and generators \(\Gamma_i(\cdot)\)
\State Initialize \(t\gets0\) and pairwise disjoint safe sets \(S_i^\ast(0)\), \(i\in\mathcal I\)
\While{termination criterion not met}
    \ForAll{\(i\in\mathcal I\) \textbf{(in parallel)}}
        \State Measure \(x_i(t)\), observe \(\{x_j(t)\}_{j\neq i}\), and reconstruct \(\{S_j^\ast(t)\}_{j\in\mathcal I}\)
        \State Solve the local FHOCP~\eqref{prob:local_fhocp}
        \State \(u_i(t)\gets u^{\mathrm n,\ast}_{i,(0|t)}=u^{\mathrm c,\ast}_{i,(0|t)}\)
    \EndFor
    \State Apply \(\{u_i(t)\}_{i\in\mathcal I}\) and measure \(\{x_i(t^+)\}_{i\in\mathcal I}\)
    \ForAll{\(i\in\mathcal I\)}
        \State Compute \(\tilde S_i(t^+)\) via~\eqref{eq:candidate_safe_set}
        \State Update \(S_i^\ast(t^+)\) via~\eqref{eq:freeze_shift_update} and \(\hat J_i^{\mathrm c}(t^+)\) via~\eqref{eq:lyap_bound_update_main}
    \EndFor
    \State \(t\gets t^+\)
\EndWhile
\end{algorithmic}
\end{algorithm}
\subsection{Plug-and-play (\textit{PnP}) operation}
\label{subsec:plug_and_play}
\textit{PnP} operation allows agents to enter or leave the workspace online while preserving safety and recursive feasibility without centralized coordination or access to past communication logs. At time \(t_{\mathrm{join}}\), a joining agent observes the current states
of the already active agents and evaluates the state-induced safe sets
\(\Gamma_j(x_j(t_{\mathrm{join}}))\). If two such sets overlap, the
corresponding candidates must have been rejected by the preceding
\textit{FoS} update, i.e.,
\(\chi_j(t_{\mathrm{join}}^-)=1\) for the involved agents. Hence, their
current active safe sets cannot be inferred from
\(\Gamma_j(x_j(t_{\mathrm{join}}))\) alone and are replaced by the
history-free overapproximation. If no such rejection is indicated, the active safe-sets of other agents can be generated by \(\Gamma_j\).


\subsubsection{Frozen-set overapproximation}

For each agent $j\in\mathcal I$, let $S_j^\ast(t)$ denote the active local safe set and let $r_j>0$ denote the agent radius.
By the safe-set containment constraint~\eqref{prob:local_fhocp_containment} enforced in the local FHOCP, the current body set satisfies $\mathcal B_j(t)\subseteq S_j^\ast(t),$
which implies
\begin{equation}
\|p_j(t)-c_j(t)\|\le R_j(t)-r_j.
\label{eq:pp_center_distance_bound}
\end{equation}
Assume moreover that a known upper bound $R_{j,\max}>0$ on the active safe-set radius is available, i.e.,
\begin{equation}
R_j(t)\le R_{j,\max},
\qquad \forall t\in\mathbb Z_+.
\label{eq:pp_radius_upper_bound}
\end{equation}
Such a bound may be obtained from the chosen safe-set generator and the admissible state set.

\begin{definition}[History-free reconstruction ball]
\label{def:pp_reconstruction_ball_frozen}
Given the current position $p_j(t)\in\mathbb R^{n_p}$ of agent $j$ and an upper bound $R_{j,\max}$ on the possibly frozen safe-set radius, define the \emph{history-free reconstruction ball}
\begin{equation}
S_j^{\mathrm{rec}}(t)
:=
\mathbb{B}\!\left(p_j(t),\,R_j^{\mathrm{rec}}\right),
\quad
R_j^{\mathrm{rec}}:=2R_{j,\max}-r_j.
\label{eq:pp_rec_ball}
\end{equation}
\end{definition}

\begin{lemma}[Outer approximation of frozen sets]
\label{lem:pp_outer_approx_frozen}
Suppose that \eqref{eq:pp_center_distance_bound} and \eqref{eq:pp_radius_upper_bound} hold.
Then, for any time $t$ at which agent $j$ is frozen, the active safe set is contained in the reconstruction ball, i.e.,
\begin{equation}
S_j^\ast(t)\subseteq S_j^{\mathrm{rec}}(t).
\label{eq:pp_outer_inclusion}
\end{equation}
\end{lemma}
\begin{proof}
Fix such a time $t$ and let $q\in S_j^\ast(t)=\mathbb B(c_j(t),R_j(t))$.
By the triangle inequality,
\[
\|q-p_j(t)\|
\le
\|q-c_j(t)\|+\|c_j(t)-p_j(t)\|.
\]
The first term satisfies $\|q-c_j(t)\|\le R_j(t)$ by definition of $S_j^\ast(t)$, and the second term satisfies $\|c_j(t)-p_j(t)\|\le R_j(t)-r_j$
by \eqref{eq:pp_center_distance_bound}.
Hence,
\[
\|q-p_j(t)\|
\le
2R_j(t)-r_j
\le
2R_{j,\max}-r_j
=
R_j^{\mathrm{rec}}.
\]
Therefore,
\[
q\in \mathbb B\bigl(p_j(t),R_j^{\mathrm{rec}}\bigr)=S_j^{\mathrm{rec}}(t),
\]
which proves \eqref{eq:pp_outer_inclusion}.
\end{proof}

Lemma~\ref{lem:pp_outer_approx_frozen} provides a conservative measurement-based substitute for an unknown frozen safe set. For non-frozen agents, no overapproximation is needed because the active safe set is reconstructable from the current state via \(\Gamma_j\).
\begin{assumption}[Feasible join initialization]
\label{ass:pp_join_feasible}
At time $t_{\mathrm{join}}$, a joining agent $i$ can select an initial active safe set
\(
S_i^\ast(t_{\mathrm{join}})
\)
according to the deterministic safe-set construction framework such that it does not overlap with the safe-set representations of already active agents, i.e.,
\begin{equation}
S_i^\ast(t_{\mathrm{join}})\cap \hat S_j(t_{\mathrm{join}})=\emptyset,
\qquad \forall j\in\mathcal I,
\label{eq:pp_join_nonoverlap}
\end{equation}
where
\[
\hat S_j(t_{\mathrm{join}})
:=
\begin{cases}
S_j^\ast(t_{\mathrm{join}}), & \text{if } \chi_j(t_{\mathrm{join}})=0,\\[0.2em]
S_j^{\mathrm{rec}}(t_{\mathrm{join}}), & \text{if } \chi_j(t_{\mathrm{join}})=1.
\end{cases}
\]
\end{assumption}

\subsubsection{Join and leave protocol}
\begin{algorithm}[H]
\caption{\emph{PnP} join protocol}
\label{alg:join_protocol}
\begin{algorithmic}[1]
\Require Safe-set generators \(\Gamma_j(\cdot)\), radius bounds \(R_{j,\max}\), and local FHOCP ingredients
\State Obtain \(\{x_j(t_{\mathrm{join}})\}_{j\in\mathcal I}\) and infer \(\{\chi_j(t_{\mathrm{join}}^-)\}_{j\in\mathcal I}\) for overlapping agents from \(\Gamma_j(x_j(t_{\mathrm{join}}))\)
\State Construct the safe-set representations $\forall j\in\mathcal I$
\[
\hat S_j(t_{\mathrm{join}})
=
\left\{
\begin{array}{@{}ll@{}}
S_j^{\mathrm{rec}}(t_{\mathrm{join}}),
&\!\! \chi_j(t_{\mathrm{join}}^-)=1.\\[0.2em]
\Gamma_j(x_j(t_{\mathrm{join}})),
&\!\! \text{otherwise}.
\end{array}
\right.
\]
\State Choose \(S_i^\ast(t_{\mathrm{join}})=\Gamma_i(x_i(t_{\mathrm{join}}))\) such that
\[
S_i^\ast(t_{\mathrm{join}})\cap \hat S_j(t_{\mathrm{join}})=\emptyset,
\qquad \forall j\in\mathcal I.
\]
\If{no such \(S_i^\ast(t_{\mathrm{join}})\) exists}
    \State Reject or postpone the join
    \State \textbf{return}
\EndIf
\State Solve the FHOCP~\eqref{prob:local_fhocp} for agent \(i\) using \(S_i^\ast(t_{\mathrm{join}})\)
\State Apply the shared first input and update \(\mathcal I\gets \mathcal I\cup\{i\}\)
\end{algorithmic}
\end{algorithm}
A leaving protocol is immediate in the decentralized setting: When an agent exits, it is removed from the set of observed agents and from the locally reconstructed constraints. The remaining agents continue to solve their local contingency MPC problems, and the closed-loop guarantees are preserved by construction.

% ============================================================================
\section{System-theoretic analysis}
\label{sec:analysis}
% ============================================================================
This section establishes the main closed-loop guarantees of the proposed decentralized contingency MPC scheme combined with the \textit{FoS} update rule.
First, for a fixed set of active agents, recursive feasibility and collision avoidance are addressed.
After that, the Lyapunov-type decrease induced by the shifted-tail bound update is analyzed and the corresponding convergence statement is derived.
The \emph{PnP} protocol introduced in Section~5 is treated separately and does not enter the core recursive-feasibility statement below.
The recursive-feasibility and Lyapunov-type arguments developed in this section follow the classical shifted-candidate logic of MPC and NMPC, adapted here to the decentralized contingency formulation with \textit{FoS}-based safe-set updates~\cite{ChenAllgoewer1998,Mayne2000}.
The detailed proofs are deferred to the Appendix.
%\subsection{Recursive Feasibility and Collision Avoidance}
\begin{assumption}[Initial feasibility and separation]
\label{ass:initial_feasibility}
At time $t_0=0$, the local FHOCP~\eqref{prob:local_fhocp} is feasible for every agent $i\in\mathcal I$, and the initial active local safe sets are pairwise disjoint, i.e.,
\[
S_i^\ast(0)\cap S_j^\ast(0)=\emptyset,
\qquad \forall i\neq j.
\]
\end{assumption}

\begin{assumption}[Nominal completion]
\label{ass:nominal_completion}
Whenever a feasible contingency candidate at time $t^+$ is available, there exists a feasible nominal candidate at time $t^+$ satisfying the nominal initialization, dynamics, and state/input constraints in~\eqref{prob:local_fhocp}.
This is, for example, satisfied if the nominal part is chosen identical to the contingency part.
\end{assumption}

\begin{lemma}[Disjointness invariance]
\label{lem:disjoint_invariance}
Under Assumption~\ref{ass:initial_feasibility} and the update rule~\eqref{eq:freeze_shift_update}, the active local safe sets remain pairwise disjoint for all $t\in\mathbb{Z}_+$, i.e.,
\[
S_i^\ast(t)\cap S_j^\ast(t)=\emptyset,
\qquad \forall i\neq j.
\]
\end{lemma}

\begin{theorem}[Recursive feasibility]
\label{thm:rec_feas}
Suppose that Assumptions~\ref{ass:exact_dynamics}--\ref{ass:contingency_recoverability}, \ref{ass:initial_feasibility}, and~\ref{ass:nominal_completion} hold.
Consider the proposed decentralized contingency MPC scheme with the \textit{FoS} update rule~\eqref{eq:freeze_shift_update} and the shifted-tail update of the Lyapunov bound~\eqref{eq:lyap_bound_update_main}.
Then the local FHOCP~\eqref{prob:local_fhocp} remains feasible for all times, i.e., for every agent $i\in\mathcal I$, \eqref{prob:local_fhocp} is feasible at every  $t\in\mathbb Z_+$.
\end{theorem}

\begin{theorem}[Closed-loop collision avoidance]
\label{thm:safety}
Suppose that the assumptions of Theorem~\ref{thm:rec_feas} hold.
Then, under the proposed decentralized contingency MPC scheme with the \textit{FoS} update rule, the closed-loop execution is collision-free for all times, i.e.,
\[
\mathcal B_i(t)\cap \mathcal B_j(t)=\emptyset,
\qquad \forall i\neq j,\ \forall t\in\mathbb Z_+.
\]
\end{theorem}

\begin{corollary}[Preservation under \emph{PnP} events]
\label{cor:pnp_preservation}
Consider the proposed decentralized contingency MPC scheme with the FoS update rule.

\emph{Join:}
Let a new agent $i_{\mathrm{new}}\notin\mathcal I$ request insertion at time $t_{\mathrm{join}}$.
If Assumption~\ref{ass:pp_join_feasible} holds, then recursive feasibility and collision avoidance are preserved for the augmented active set
\[
\mathcal I^+ := \mathcal I \cup \{i_{\mathrm{new}}\} \quad \forall t\ge t_{\mathrm{join}}.
\]
\emph{Leave:}
If an agent $j\in\mathcal I$ leaves the workspace at time $t_{\mathrm{leave}}$ and is removed from the locally reconstructed constraints of the remaining agents. Recursive feasibility and collision avoidance are preserved for the reduced active set
\[
\mathcal I^- := \mathcal I \setminus \{j\} \quad \forall t \geq t_{\mathrm{leave}}.
\]
\end{corollary}
The convergence mechanism is induced by the Lyapunov-type constraint~\eqref{prob:local_fhocp_lyap}.
The argument follows the standard receding-horizon value-function logic: feasibility of the shifted contingency candidate yields an explicit feasible upper bound for the successor problem, and the imposed Lyapunov constraint then enforces a monotone decrease of the optimal contingency-cost sequence. This is the same shifted-candidate proof architecture used in classical MPC stability proofs; see, e.g.,~\cite{ChenAllgoewer1998,Mayne2000}.

\begin{definition}[$\mathcal K_\infty$ function]
\label{def:Kinf}
A continuous function $\alpha:\mathbb{R}_+\to\mathbb{R}_+$ is said to belong to class $\mathcal K_\infty$ if
\(
\alpha(0)=0,
\)
$\alpha$ is strictly increasing, and
\(
\alpha(s)\to\infty
\ \text{as } s\to\infty.
\)
\end{definition}

\begin{assumption}[Regularity and coercivity]
\label{ass:regularity_cost}
Let
\(
\mathcal Z_i := \mathcal X_i\times\mathcal U_i,
\)
and let $\overline{\mathcal Z}_i$ denote its closure.
For each agent $i\in\mathcal I$, the following properties hold:

\begin{enumerate}
    \item \textbf{Continuity of dynamics and cost.}
    The mappings
    \(
    f_i:\mathbb R^{n_i}\times\mathbb R^{m_i}\to\mathbb R^{n_i},
    \
    \ell_i^{\mathrm c}:\mathbb R^{n_i}\times\mathbb R^{m_i}\to\mathbb R_+,
    \)
    and the offset cost
    \(
    V_i^{\mathrm c}:\mathbb R^{n_i}\times\mathbb R^{n_i}\to\mathbb R_+
    \)
    are continuous on the relevant admissible sets.

    \item \textbf{Incremental continuity of the dynamics.}
    There exists a function $\alpha_{f,i}\in\mathcal K_\infty$ such that
    \begin{align*}
    \|f_i(x,u)-f_i(\tilde x,\tilde u)\|
    \le
    \alpha_{f,i}\!\left(
    \left\|
    \begin{bmatrix}
    x-\tilde x\\
    u-\tilde u
    \end{bmatrix}
    \right\|
    \right),
    \\
    \forall (x,u),(\tilde x,\tilde u)\in\overline{\mathcal Z}_i.
    \end{align*}

    \item \textbf{Positive definiteness.}
    There exist functions
    \(
    \underline{\alpha}_{\ell,i},\,\overline{\alpha}_{\ell,i}\in\mathcal K_\infty
    \)
    such that
    \begin{equation}
    \begin{aligned}
    \underline{\alpha}_{\ell,i}\!\bigl(\|(e,\delta u)\|\bigr)
    \le
    \ell_i^{\mathrm c}(e,\delta u)
    \le
    \overline{\alpha}_{\ell,i}\!\bigl(\|(e,\delta u)\|\bigr),
    \\
    \forall (e,\delta u)\in\overline{\mathcal Z}^{\mathrm e}_i,
    \label{eq:stage_cost_bounds}
    \end{aligned}
    \end{equation}
    where $\overline{\mathcal Z}^{\mathrm e}_i$ denotes the relevant closed admissible set in error coordinates: $e = x_i^c-\bar{x}_i^c$ and $\delta u = u_i^\mathrm{c}-\bar{u}_i^c$.
    In particular, $\ell_i^{\mathrm c}$ is positive definite with respect to $(e,\delta u)=(0,0)$.

    \item \textbf{Positive definiteness of the offset cost.}
    There exist functions
    \(
    \underline{\alpha}_{V,i},\,\overline{\alpha}_{V,i}\in\mathcal K_\infty
    \)
    such that
    \[
    \underline{\alpha}_{V,i}\!\bigl(\|\bar{x}_i^c- x_i^{\mathrm{ref}}\|\bigr)
    \le
    V_i^{\mathrm c}(\bar{x}_i^c,x_i^{\mathrm{ref}})
    \le
    \overline{\alpha}_{V,i}\!\bigl(\|\bar{x}_i^c- x_i^{\mathrm{ref}}\|\bigr)
    \]
    for all relevant admissible equilibrium candidates $\bar x_i^{\mathrm c}$.
\end{enumerate}
\end{assumption}
For a fixed agent \(i\in\mathcal I\), define the equilibrium-tracking
errors associated with the optimal contingency equilibrium by
\[
e_i(t):=x_i(t)-\bar x_i^{\mathrm c,\ast}(t),
\qquad
\delta u_i(t):=u_i(t)-\bar u_i^{\mathrm c,\ast}(t),
\]
where
\(
u_i(t)=u_{i,(0|t)}^{\mathrm c,\ast}
      =u_{i,(0|t)}^{\mathrm n,\ast}
\)
is the applied shared first input and $x_i(t)$ is as in (\ref{prob:local_fhocp_init}).
\begin{proposition}%[Lyapunov-type decrease]
\label{prop:lyap_decrease}
Fix an agent \(i\in\mathcal I\) and suppose that recursive feasibility of
the local FHOCP~\eqref{prob:local_fhocp} holds for all
\(t\in\mathbb Z_+\). Assume that the Lyapunov
constraint~\eqref{prob:local_fhocp_lyap} is enforced at every time step,
that the bound \(\hat J_i^{\mathrm c}(\cdot)\) is updated by the
shifted-tail rule~\eqref{eq:lyap_bound_update_main}, and that
Assumption~\ref{ass:regularity_cost} holds. Then
\begin{equation}
J_i^{\mathrm c,\ast}(t^+)
\le
J_i^{\mathrm c,\ast}(t)
-
\ell_i^{\mathrm c}\!\bigl(e_i(t),\delta u_i(t)\bigr),
\qquad \forall t\in\mathbb Z_+.
\label{eq:lyap_decrease_main}
\end{equation}
In particular, the sequence $\{J_i^{\mathrm c,\ast}(t)\}_{t\in\mathbb Z_+}$ is monotonically nonincreasing and lower bounded, hence convergent, and
\begin{equation}
\sum_{t=0}^{\infty}
\ell_i^{\mathrm c}\!\bigl(e_i(t),\delta u_i(t)\bigr)
<\infty.
\label{eq:summable_stage_main}
\end{equation}
Consequently,
\(
\ell_i^{\mathrm c}\!\bigl(e_i(t),\delta u_i(t)\bigr)\to 0
\ \text{as } t\to\infty,
\)
and by the lower bound in~\eqref{eq:stage_cost_bounds},
\(
\|(e_i(t),\delta u_i(t))\|\to 0
\ \text{as } t\to\infty.
\)
\end{proposition}
\begin{remark}[Interpretation of Proposition~\ref{prop:lyap_decrease}]
\label{rem:prop_lyap_interpretation}
Proposition~\ref{prop:lyap_decrease} is the precise Lyapunov-type statement available without any additional assumption on the time evolution of the selected contingency equilibrium.
It guarantees monotone decrease of the optimal contingency cost, summability of the stage cost, and asymptotic convergence of the closed-loop trajectory in the equilibrium-tracking coordinates $(e_i(t),\delta u_i(t))$.
\end{remark}
\begin{corollary}[Convergence to a fixed equilibrium]
\label{cor:fixed_ss_convergence}
Suppose that the assumptions of Proposition~\ref{prop:lyap_decrease} hold and that there exists \(\bar t\in\mathbb Z_+\) such that, for all \(t\ge\bar t\), \(\bar x_i^{\mathrm c,\ast}(t)=\bar x_i^{\mathrm c,\star}\) and \(\bar u_i^{\mathrm c,\ast}(t)=\bar u_i^{\mathrm c,\star}\). Then \(x_i(t)\to \bar x_i^{\mathrm c,\star}\) and \(u_i(t)\to \bar u_i^{\mathrm c,\star}\) as \(t\to\infty\).
\end{corollary}
% ============================================================================
\section{Simulation Study}
\label{sec:sim}
% ============================================================================
The simulation results presented in the following are illustrated in the supplementary video available at \url{https://youtu.be/uBLV58GRvFw}.
\subsection{Setup}
\label{subsec:sim_setup}

The general framework developed in the previous sections is now instantiated for a specific agent model in order to assess the practical performance of the proposed decentralized contingency MPC scheme.
For the simulation study, each agent is modeled by a discrete-time double integrator in a bounded planar workspace
\(
\mathcal W \subset \mathbb R^2, \ \text{i.e.,} \ n_p = 2.
\)
For each agent $i\in\mathcal I$, the state is given by
\(
x_i(t)=
[
p_i(t)\ \
v_i(t)
]^\top
\in\mathbb R^4,
\ 
p_i(t)\in\mathbb R^2,
\
v_i(t)\in\mathbb R^2,
\)
and the input is
\(
u_i(t)\in\mathbb R^2.
\)
The discrete-time dynamics are
\begin{equation}
x_i(t^+)=f(x_i(t),u_i(t)) = Ax_i(t)+Bu_i(t),
\label{eq:sim_double_integrator}
\end{equation}
with sampling time $T_s=0.1\,\mathrm{s}$ and system matrices
\[
A=
\begin{bmatrix}
I_2 & T_s I_2\\
0_2 & I_2
\end{bmatrix},
\qquad
B=
\begin{bmatrix}
\frac{1}{2}T_s^2 I_2\\
T_s I_2
\end{bmatrix},
\]
where $I_2\in\mathbb R^{2\times 2}$ denotes the identity matrix and $0_2\in\mathbb R^{2\times 2}$ the zero matrix.
The position output is
\(
p_i(t)=C_i x_i(t),
\
C_i=[ I_2 \ \  0_2 ].
\)
Each agent body is approximated by a disc of radius
\(
r_i=0.2\,\mathrm{m}.
\)
State and input bounds are imposed through the velocity and acceleration limits
\(
\|v_i(t)\|\le 3\,\mathrm{m/s},
\
\|u_i(t)\|\le 3.5\,\mathrm{m/s}^2.
\)
The nominal prediction horizon is chosen as \(N_{\mathrm n}=20\). For the
double-integrator model, the contingency horizon \(N_{\mathrm c}\) is
chosen to allow one shared first input followed by admissible braking to
a stopped equilibrium,
\(
\bar x_i^{\mathrm c} =
[
\bar p_i^{\mathrm c} \ \ 
0
]^\top,
\ 
\bar u_i^{\mathrm c}=0.
\)
A sufficient horizon length is
\begin{equation}
N_{\mathrm c}
\ge
1+\left\lceil\frac{v_{\max}}{a_{\max}T_s}\right\rceil .
\label{eq:sim_Nc_choice}
\end{equation}
For the double-integrator model, the active safe-set center in
\eqref{eq:local_safe_set} is chosen as the current position,
\(c_i(t)=p_i(t)\). The corresponding stopping radius is
\[
R_i(t)
=
\frac{1}{2}T_s^2\|u\|
+
T_s\|v\|
+
\frac{\|v+T_su\|^2}{2|a_{\mathrm{min}}|}
+
r_i .
\]
The first two terms bound the displacement during the shared first input
\(u\in\mathcal U\), while the third term bounds the subsequent braking
distance. This particular radius is not essential; any choice satisfying
Assumption~\ref{ass:contingency_recoverability} is admissible.

\subsubsection{MPC Implementation}
All local optimal control problems are implemented in \textsc{CasADi} and solved using \textsc{IPOPT}~\cite{IPOPT}.
Solver tolerances are set to $10^{-5}$ for feasibility and optimality, and the maximum number of iterations is capped at $2000$.
Warm-starting is used throughout.
\label{subsec:mpc_impl}
\subsubsection*{Tracking and equilibrium optimization}
The quadratic objective in the form of 
\begin{align*}
&J_i\big(X_i^{\mathrm{n}}(t), U_i^{\mathrm{n}}(t), \bar{x}^{\mathrm{c}}_i(t), x_i^{\mathrm{ref}}\big)= \\ 
&\sum_{k=0}^{N_{\mathrm n}-1}
\Big(
\|p^{\mathrm n}_{i,(k|t)}-p_i^{\mathrm{ref}}\|_{Q_p}^2
+\|v^{\mathrm n}_{i,(k|t)}\|_{Q_v}^2
+\|u^{\mathrm n}_{i,(k|t)}\|_{R_u}^2
\Big) \\
& \qquad \ +\|p^{\mathrm n}_{i,(N_{\mathrm n}|t)}-p_i^{\mathrm{ref}}\|_{Q_p}^2 +\gamma\,V_i^{\mathrm{c}}\!\big(\bar{x}^{\mathrm{c}}_i(t), x_i^{\mathrm{ref}}\big),
\end{align*}
with diagonal weights
\(
Q_p = Q_v = R_u=I_2,
\)
and
\[
V_i^{\mathrm{c}}\!\big(\bar{x}^{\mathrm{c}}_i(t), x_i^{\mathrm{ref}}\big) =\|\bar x_i^{\mathrm c}(t)-x_i^{\mathrm{ref}}\|_{P_s}^2,\ \
x^{\mathrm{ref}}_i=\begin{bmatrix}p_i^{\mathrm{ref}}&0\end{bmatrix}^\top,
\]
where $P_s = I_4$ and $\gamma=0.1$ is used in the simulations.

\subsubsection*{Contingency cost}
The Lyapunov bound uses the safe-cost functional
\begin{align*}
J_i^{\mathrm c}(t)
=
\sum_{k=0}^{N_{\mathrm c}-1}&
\Big(
\|x^{\mathrm c}_{i,(k|t)}-\bar x_i^{\mathrm c}(t)\|_{Q_s}^2
+\|u^{\mathrm c}_{i,(k|t)}\|_{R_s}^2
\Big) \\
&+\|\bar x_i^{\mathrm c}(t)-x_i^{\mathrm{ref}}\|_{P_s}^2,
\end{align*}
implemented with $Q_s=0.1\cdot I_4$, $R_s= 0.1 \cdot I_2$, and $P_s$ as before.

\subsubsection*{Nominal soft collision-avoidance constraints}
For each other agent $j\neq i$ and each nominal stage $k\in\mathbb{Z}_0^{N_{\mathrm n}-1}$ we introduce nonnegative slacks
$s_{ij,k}\ge 0$, and impose the squared-distance constraint
\begin{equation}
\|p^{\mathrm n}_{i,(k|t)}-c_j(t)\|^2 + s_{ij,k}
\;\ge\;
d_{ij,k}^2,
\label{eq:soft_dist_gate}
\end{equation}
where $c_j(t)$ and $R_j(t)$ denote the \emph{currently reconstructed} active safe-set center and radius of agent $j$, and
\[
d_{ij,k}
:=
R_{\mathrm{s}}\big(x^{\mathrm n}_{i,(k|t)}\big) + R_j(t).
\]
The slack variables are penalized in the nominal objective by
\[
J_i \leftarrow J_i
+
\rho_{\mathrm{nom}}
\sum_{j\in\mathcal I\setminus\{i\}}
\sum_{k=0}^{N_{\mathrm n}-1}
s_{ij,k}.
\]
with $\rho_{\mathrm{nom}}=100$. 
\subsubsection{Increasing agent density}
This scenario class evaluates scalability with respect to agent density. The increasing-density study uses \(M=5,10,20\) agents with randomly
sampled initial and reference positions under a minimum separation. Fig.~\ref{fig:inc_density_traj} shows collision-free
closed-loop trajectories for all densities. The normalized contingency
cost in Fig.~\ref{fig:inc_density_cost} decreases as enforced by
\eqref{prob:local_fhocp_lyap}; Higher densities mainly increase the
transient phase due to stronger local interactions.
\begin{figure}[h]
    \centering
    % upper row
    \begin{minipage}[t]{0.5\linewidth}
        \centering
        \includegraphics[width=\linewidth]{Figures/5_Paper.png}
    \end{minipage}\hfill
    \begin{minipage}[t]{0.5\linewidth}
        \centering
        \includegraphics[width=\linewidth]{Figures/10_Paper.png}
    \end{minipage}
    \vspace{0pt}
    % lower row
    \begin{minipage}[t]{0.5\linewidth}
        \centering
        \includegraphics[width=\linewidth]{Figures/20_Paper.png}
    \end{minipage}
    \caption{Representative closed-loop trajectories in the increasing-density study: $M=5$ (upper left), $M=10$ (upper right), and $M=20$ (bottom). Filled circles denote initial positions and star markers denote reference positions. The gray objects represent static obstacles. The trajectories remain collision-free in all cases.}
    \label{fig:inc_density_traj}
\end{figure}
\begin{figure}[h!]
        \includegraphics[width=1\linewidth]{Figures/Lyaponov_10_Paper.png}
    \caption{Normalized contingency-cost evolution $J_i^c(t)/J_i^c(0)$ for $M = 10$ agents (see Fig. \ref{fig:inc_density_traj}). The trajectories satisfy the Lyapunov-decreasing behavior induced by the shifted-tail contingency-cost update (constraint \eqref{prob:local_fhocp_lyap}).}
    \label{fig:inc_density_cost}
\end{figure}
\subsubsection{Bottleneck}
The bottleneck scenario stresses conflict resolution in a narrow passage
formed by two static obstacles. Fig.~\ref{fig:FieldBottleneck} shows
that all agents pass the constriction without collisions. The distance
signals and freeze count in Figs.~\ref{fig:AgentDistanceBottleneck}--\ref{fig:FreezeCountBottleneck}
show that close encounters near the passage trigger temporary freeze
events, which prevent unsafe safe-set updates while preserving
closed-loop progress.
\begin{figure}[h]
    \centering
    \includegraphics[width=0.6\linewidth]{Figures/Bottleneck_Field.png}
    \caption{Closed-loop trajectories in the bottleneck scenario. Filled circles denote initial positions and star markers denote target positions. The gray obstacles create a narrow passage that induces strong local interactions.}
    \label{fig:FieldBottleneck}
\end{figure}
\begin{figure}[h]
    \centering
    \includegraphics[width=0.9\linewidth]{Figures/AgentDistBottleneck.png}\\[-1.2em]
    \includegraphics[width=0.9\linewidth]{Figures/SafeSetDistBottleneck.png}
    \caption{Distance evolution in the \emph{bottleneck} scenario. Top: pairwise distances between agents; the dashed line indicates the minimum admissible distance. Bottom: distances between agent safe sets; values close to zero indicate strong interaction and freeze-operator activation.}
    \label{fig:AgentDistanceBottleneck}
\end{figure}
\begin{figure}[h]
    \centering
    \includegraphics[width=0.9\linewidth]{Figures/FreezeActivationBottleneck.png}
    \caption{Number of frozen agents over time in the \emph{bottleneck} scenario. The freeze events are sparse and temporary.}%, indicating that local conflicts are resolved after short interaction phases.}
    \label{fig:FreezeCountBottleneck}
\end{figure}
\subsubsection{PnP}
The \emph{PnP} experiment evaluates the proposed scheme under online changes in the set of active agents. Starting from a feasible multi-agent configuration, three additional agents join the workspace during closed-loop operation and are initialized according to the join protocol in Section~\ref{subsec:plug_and_play}. Later, a different agent leaves the workspace and is removed from the locally reconstructed constraints of the remaining agents.

Fig.~\ref{fig:PlugAndPlayFigure} summarizes the experiment. The results illustrate that insertion and removal can be handled without trajectory communication and without modifying the local MPC problem structure of the remaining agents. In particular, the deterministic \textit{FoS} update and the (history-free) reconstruction at the join time allow the scheme to preserve collision avoidance and feasibility despite the changing agent configuration.

\begin{figure}[h!]
    \centering
    \includegraphics[width=0.9\linewidth]{Figures/PlugAndPlay.png}
    \caption{\emph{PnP} scenario with online agent insertion and removal. \textbf{Left:} Closed-loop trajectories in a cluttered workspace (gray obstacles). Filled circles denote initial positions and stars denote final positions; colors correspond to individual agents. \textbf{Right:} Activity timeline indicating when each agent is active/inactive; vertical dashed lines mark join/leave events. The proposed scheme maintains safe, collision-free motion throughout despite the changing set of active agents.}
    \label{fig:PlugAndPlayFigure}
\end{figure}
% ============================================================================
\section{Conclusions and outlook}
% ============================================================================
The proposed decentralized contingency MPC scheme provides a provably safe mechanism for local interaction in strongly coupled multi-agent systems without trajectory exchange or communication. Safety is ensured through local constraint reconstruction and the deterministic \emph{FoS} update of the active safe sets. Recursive feasibility, collision-free closed-loop execution, and Lyapunov-type convergence follow from the shifted-tail construction and the associated decrease condition. The simulations indicate moderate conservatism even in dense and cluttered environments.

The safe-set principle is also aligned with intuitive behavior in traffic: a driver maintains a safety margin such that an emergency maneuver remains feasible. Similarly, the local safe sets ensure that a feasible contingency action remains available during multi-agent interaction.

Future work should address disturbances, model mismatch, sensing noise or delays in reconstructed neighbor safe sets, and heterogeneous agent dynamics and constraints. Another direction is a memoryless safe-set update, enabling active safe sets or tight outer approximations to be reconstructed from instantaneous information only. This would simplify decentralized implementation and reduce the conservatism of \emph{PnP} operation.
\appendix


% ============================================================================
% Appendix replacement (from \appendix onward)
% ============================================================================

\appendix
\numberwithin{equation}{section}

\makeatletter
\renewcommand\@seccntformat[1]{%
  \csname the#1\endcsname.\quad
}
\makeatother
\section*{Appendix}
\section{Auxiliary shifted-tail lemmas}

The recursive-feasibility and Lyapunov-type arguments both rely on the same shifted contingency candidate. To avoid duplicate arguments and cross-references between proofs, the two key shifted-tail facts are collected first.

\begin{lemma}[Feasibility of the shifted candidate]
\label{lem:appendix_shifted_candidate_feasible}
Fix an agent $i\in\mathcal I$ and a time $t\in\mathbb Z_+$ at which the local FHOCP~\eqref{prob:local_fhocp} is feasible. Let
\begin{align*}
\Big(
&\{x_{i,(k|t)}^{\mathrm n,\ast}\}_{k=0}^{N_{\mathrm n}},
\{u_{i,(k|t)}^{\mathrm n,\ast}\}_{k=0}^{N_{\mathrm n}-1},
\{x_{i,(k|t)}^{\mathrm c,\ast}\}_{k=0}^{N_{\mathrm c}},
\{u_{i,(k|t)}^{\mathrm c,\ast}\}_{k=0}^{N_{\mathrm c}-1},\\
&\qquad \qquad \qquad \bar x_i^{\mathrm c,\ast}(t),
\bar u_i^{\mathrm c,\ast}(t)
\Big)
\end{align*}
be an optimal feasible solution. Define, at time $t^+$,
\begin{align}
\tilde x_{i,(k|t^+)}^{\mathrm c}
&:=x_{i,(k+1|t)}^{\mathrm c,\ast},
\qquad k=0,\dots,N_{\mathrm c}-1,
\label{eq:app_shift_x}
\\
\tilde u_{i,(k|t^+)}^{\mathrm c}
&:=u_{i,(k+1|t)}^{\mathrm c,\ast},
\qquad k=0,\dots,N_{\mathrm c}-2,
\label{eq:app_shift_u}
\\
\tilde u_{i,(N_{\mathrm c}-1|t^+)}^{\mathrm c}
&:=\bar u_i^{\mathrm c,\ast}(t),
\label{eq:app_shift_u_terminal}
\\
\tilde{\bar x}_i^{\mathrm c}(t^+)
&:=\bar x_i^{\mathrm c,\ast}(t),
\qquad
\tilde{\bar u}_i^{\mathrm c}(t^+):=\bar u_i^{\mathrm c,\ast}(t).
\label{eq:app_shift_bar}
\end{align}
Then the candidate
\[
\Big(
\{\tilde x_{i,(k|t^+)}^{\mathrm c}\}_{k=0}^{N_{\mathrm c}},
\{\tilde u_{i,(k|t^+)}^{\mathrm c}\}_{k=0}^{N_{\mathrm c}-1},
\tilde{\bar x}_i^{\mathrm c}(t^+),
\tilde{\bar u}_i^{\mathrm c}(t^+)
\Big)
\]
satisfies the contingency-part constraints
\eqref{prob:local_fhocp_init},
\eqref{prob:local_fhocp_dyn_c},
\eqref{prob:local_fhocp_bounds_c},
\eqref{prob:static_obstacle_constraintsMPC},
\eqref{prob:local_fhocp_terminal_state}--\eqref{prob:local_fhocp_terminal_position},
\eqref{prob:local_fhocp_containment},
and \eqref{prob:local_fhocp_tail} at time $t^+$.
\end{lemma}

\begin{proof}
By feasibility at time $t$,
\begin{equation}
x_{i,(N_{\mathrm c}|t)}^{\mathrm c,\ast}
=
\bar x_i^{\mathrm c,\ast}(t),
\
\bar x_i^{\mathrm c,\ast}(t)
=
f_i\!\bigl(\bar x_i^{\mathrm c,\ast}(t),\bar u_i^{\mathrm c,\ast}(t)\bigr).
\label{eq:app_terminal_eq}
\end{equation}
The applied input is
\(
u_i(t)=u_{i,(0|t)}^{\mathrm n,\ast}=u_{i,(0|t)}^{\mathrm c,\ast}.
\)
Hence, by Assumption~\ref{ass:exact_dynamics},
\begin{equation}
x_i(t^+)
=
f_i\!\bigl(x_{i,(0|t)}^{\mathrm c,\ast},u_{i,(0|t)}^{\mathrm c,\ast}\bigr)
=
x_{i,(1|t)}^{\mathrm c,\ast}.
\label{eq:app_successor_state}
\end{equation}
We first define the terminal candidate state by
\begin{equation}
\tilde x_{i,(N_{\mathrm c}|t^+)}^{\mathrm c}
:=
f_i\!\bigl(
\tilde x_{i,(N_{\mathrm c}-1|t^+)}^{\mathrm c},
\tilde u_{i,(N_{\mathrm c}-1|t^+)}^{\mathrm c}
\bigr).
\label{eq:app_shift_terminal_state}
\end{equation}
Using \eqref{eq:app_shift_x}, \eqref{eq:app_shift_u_terminal}, and \eqref{eq:app_terminal_eq},
\begin{equation}
\begin{aligned}
\tilde x_{i,(N_{\mathrm c}|t^+)}^{\mathrm c}
&=
f_i\bigl(x_{i,(N_{\mathrm c}|t)}^{\mathrm c,\ast},\bar u_i^{\mathrm c,\ast}(t)\bigr)\\
&=
f_i\bigl(\bar x_i^{\mathrm c,\ast}(t),\bar u_i^{\mathrm c,\ast}(t)\bigr)\\
&=
\bar x_i^{\mathrm c,\ast}(t)
=
\tilde{\bar x}_i^{\mathrm c}(t^+).
\label{eq:app_terminal_state_equals_bar}
\end{aligned}
\end{equation}
\smallskip
\noindent\textbf{Initialization and dynamics.}
By \eqref{eq:app_successor_state} and \eqref{eq:app_shift_x},
\[
\tilde x_{i,(0|t^+)}^{\mathrm c}
=
x_{i,(1|t)}^{\mathrm c,\ast}
=
x_i(t^+),
\]
so \eqref{prob:local_fhocp_init} holds for the contingency part.
For $k=0,\dots,N_{\mathrm c}-2$, feasibility at time $t$ gives
\begin{align*}
\tilde x_{i,(k+1|t^+)}^{\mathrm c}
=
x_{i,(k+2|t)}^{\mathrm c,\ast}
&=
f_i\!\bigl(x_{i,(k+1|t)}^{\mathrm c,\ast},u_{i,(k+1|t)}^{\mathrm c,\ast}\bigr)
\\
&=
f_i\!\bigl(\tilde x_{i,(k|t^+)}^{\mathrm c},\tilde u_{i,(k|t^+)}^{\mathrm c}\bigr),
\end{align*}
Thus \eqref{prob:local_fhocp_dyn_c} holds for $k=0,\dots,N_{\mathrm c}-2$, and the case $k=N_{\mathrm c}-1$ is given by \eqref{eq:app_shift_terminal_state}.

\smallskip
\noindent\textbf{State, input, and static-obstacle constraints.}
For $k=0,\dots,N_{\mathrm c}-2$,
\[
\tilde x_{i,(k|t^+)}^{\mathrm c}
=
x_{i,(k+1|t)}^{\mathrm c,\ast}\in\mathcal X_i,
\qquad
\tilde u_{i,(k|t^+)}^{\mathrm c}
=
u_{i,(k+1|t)}^{\mathrm c,\ast}\in\mathcal U_i
\]
by feasibility at time $t$.
For the last stage,
\[
\tilde x_{i,(N_{\mathrm c}-1|t^+)}^{\mathrm c}
=
x_{i,(N_{\mathrm c}|t)}^{\mathrm c,\ast}
=
\bar x_i^{\mathrm c,\ast}(t)\in\mathcal X_i,
\]
and
\[
\tilde u_{i,(N_{\mathrm c}-1|t^+)}^{\mathrm c}
=
\bar u_i^{\mathrm c,\ast}(t)\in\mathcal U_i.
\]
The same index shift shows that
\[
h\bigl(\tilde p_{i,(k|t^+)}^{\mathrm c}\bigr)\ge 0,
\qquad \forall k\in\mathbb Z_0^{N_{\mathrm c}},
\]
where $\tilde p_{i,(k|t^+)}^{\mathrm c}:=C_i\tilde x_{i,(k|t^+)}^{\mathrm c}$.

\smallskip
\noindent\textbf{Terminal constraints.}
Equation \eqref{eq:app_terminal_state_equals_bar} yields
\[
\tilde x_{i,(N_{\mathrm c}|t^+)}^{\mathrm c}
=
\tilde{\bar x}_i^{\mathrm c}(t^+),
\]
which is \eqref{prob:local_fhocp_terminal_state}.
Moreover, by \eqref{eq:app_shift_bar} and \eqref{eq:app_terminal_eq},
\[
\tilde{\bar x}_i^{\mathrm c}(t^+)
=
f_i\!\bigl(\tilde{\bar x}_i^{\mathrm c}(t^+),\tilde{\bar u}_i^{\mathrm c}(t^+)\bigr),
\]
so \eqref{prob:local_fhocp_terminal_equilibrium} holds.
Finally, feasibility at time $t$ implies
\[
C_i\bar x_i^{\mathrm c,\ast}(t)\in S_i^\ast(t).
\]
If $\chi_i(t)=1$, then $S_i^\ast(t^+)=S_i^\ast(t)$ by \eqref{eq:freeze_shift_update}, so
\[
C_i\tilde{\bar x}_i^{\mathrm c}(t^+)\in S_i^\ast(t^+).
\]
If $\chi_i(t)=0$, then $S_i^\ast(t^+)=\tilde S_i(t^+)=\Gamma_i(x_i(t^+))$, and the tail-containment argument below with $k=1$ already implies
\[
C_i\tilde{\bar x}_i^{\mathrm c}(t^+)\in S_i^\ast(t^+).
\]
Hence \eqref{prob:local_fhocp_terminal_position} holds.

\smallskip
\noindent\textbf{Active safe-set containment.}
Active safe-set containment~\eqref{prob:local_fhocp_containment}.
It remains to show that the shifted contingency candidate satisfies the active safe-set containment constraint at time~$t^+$.
We distinguish two cases.

\emph{Case 1:} $\chi_i(t)=1$.
Then, by the FoS update rule~\eqref{eq:freeze_shift_update}, the active safe set is frozen, i.e.,
\(
S_i^\ast(t^+) = S_i^\ast(t).
\)
Since the contingency trajectory at time~$t$ was feasible and satisfied~\eqref{prob:local_fhocp_containment} in~$S_i^\ast(t)$, the shifted tail remains contained in the same set.
Hence,~\eqref{prob:local_fhocp_containment} holds at time~$t^+$.

\emph{Case 2:} $\chi_i(t)=0$.
Then the active safe set is updated to the state-induced candidate set,
\(
S_i^\ast(t^+) = \Gamma_i\!\bigl(x_i(t^+)\bigr).
\)
By exact dynamics and the shared first input,
\[
x_i(t^+) = x_{i,(1|t)}^{c,\ast}.
\]
Therefore, the new active safe set is exactly the one induced by the first predicted contingency successor state.
Since the feasible solution at time~$t$ satisfies the tail-containment constraint~\eqref{prob:local_fhocp_tail}, the entire shifted contingency tail is contained in that induced safe set.
Hence,~\eqref{prob:local_fhocp_containment} also holds at time~$t^+$.

\emph{Tail-containment constraint~\eqref{prob:local_fhocp_tail}.}
Let arbitrary
\(
k\in\mathbb{Z}_{0}^{N_c}, \ l\in\mathbb{Z}_{k}^{N_c}
\)
be given for the candidate at time~$t^+$.
It must be shown that
\[
\bigl\|\tilde p^{\mathrm c}_{i,(l|t^+)}
-
c_i\!\bigl(\tilde x^{\mathrm c}_{i,(k|t^+)}\bigr)\bigr\|
\le
R_i\!\bigl(\tilde x^{\mathrm c}_{i,(k|t^+)}\bigr)-r_i.
\]
For $k\le N_c-1$ and $l\le N_c-1$, the shifted candidate satisfies
\[
\tilde x^{\mathrm c}_{i,(l|t^+)} = x^{\mathrm c,\ast}_{i,(l+1|t)},
\qquad
\tilde x^{\mathrm c}_{i,(k|t^+)} = x^{\mathrm c,\ast}_{i,(k+1|t)}.
\]
Since the safe-set generator is deterministic, this implies
\[
c_i\!\bigl(\tilde x^{\mathrm c}_{i,(k|t^+)}\bigr)
=
c_i\!\bigl(x^{\mathrm c,\ast}_{i,(k+1|t)}\bigr),
\ 
R_i\!\bigl(\tilde x^{\mathrm c}_{i,(k|t^+)}\bigr)
=
R_i\!\bigl(x^{\mathrm c,\ast}_{i,(k+1|t)}\bigr).
\]
Moreover, $l\in\mathbb{Z}_{k}^{N_c-1}$ implies $l+1\in\mathbb{Z}_{k+1}^{N_c}$.
Hence, feasibility of the solution at time~$t$ and constraint~\eqref{prob:local_fhocp_tail} yield
\[
\bigl\|p^{\mathrm c,\ast}_{i,(l+1|t)}
-
c_i\!\bigl(x^{\mathrm c,\ast}_{i,(k+1|t)}\bigr)\bigr\|
\le
R_i\!\bigl(x^{\mathrm c,\ast}_{i,(k+1|t)}\bigr)-r_i,
\]
which is exactly the desired inequality.

For $k\le N_c-1$ and $l=N_c$, we have
\[
\tilde x^{\mathrm c}_{i,(N_c|t^+)}
=
\tilde{\bar x}^{\mathrm c}_i(t^+)
=
\bar x^{\mathrm c,\ast}_i(t)
=
x^{\mathrm c,\ast}_{i,(N_c|t)}.
\]
Since now $N_c\in\mathbb{Z}_{k+1}^{N_c}$, feasibility at time~$t$ again gives
\[
\bigl\|\tilde p^{\mathrm c}_{i,(N_c|t^+)}
-
c_i\!\bigl(\tilde x^{\mathrm c}_{i,(k|t^+)}\bigr)\bigr\|
\le
R_i\!\bigl(\tilde x^{\mathrm c}_{i,(k|t^+)}\bigr)-r_i.
\]

Finally, for $k=N_c$, necessarily $l=N_c$, and the claim follows directly from
\[
\tilde x^{\mathrm c}_{i,(N_c|t^+)} = x^{\mathrm c,\ast}_{i,(N_c|t)}
\]
together with feasibility at time~$t$ for the index pair $(k,l)=(N_c,N_c)$.

Therefore, the shifted candidate satisfies
\eqref{prob:local_fhocp_tail} at time~$t^+$.
\end{proof}
\begin{lemma}[Shifted contingency-cost identity]
\label{lem:appendix_shifted_cost_identity}
Under the hypotheses of Lemma~\ref{lem:appendix_shifted_candidate_feasible}, the shifted contingency candidate defined in
\eqref{eq:app_shift_x}--\eqref{eq:app_shift_bar} satisfies
\begin{equation}
\tilde J_i^{\mathrm c}(t^+)
=
\hat J_i^{\mathrm c}(t^+),
\label{eq:app_candidate_equals_bound}
\end{equation}
where $\hat J_i^{\mathrm c}(t^+)$ is given by \eqref{eq:lyap_bound_update_main}.
\end{lemma}

\begin{proof}
By definition \eqref{eq:c_cost_def},
\begin{equation}
\begin{aligned}
\tilde J_i^{\mathrm c}(t^+)
=
\sum_{k=0}^{N_{\mathrm c}-1}
&\ell_i^{\mathrm c}\!\Big(
\tilde x_{i,(k|t^+)}^{\mathrm c}-\tilde{\bar x}_i^{\mathrm c}(t^+),
\tilde u_{i,(k|t^+)}^{\mathrm c}-\tilde{\bar u}_i^{\mathrm c}(t^+)
\Big)
\\
&\;+\;
V_i^{\mathrm c}\!\bigl(\tilde{\bar x}_i^{\mathrm c}(t^+),x_i^{\mathrm{ref}}\bigr).
\end{aligned}
\label{eq:app_shifted_candidate_cost}
\end{equation}
For $k=0,\dots,N_{\mathrm c}-2$,
\[
\tilde x_{i,(k|t^+)}^{\mathrm c}-\tilde{\bar x}_i^{\mathrm c}(t^+)
=
x_{i,(k+1|t)}^{\mathrm c,\ast}-\bar x_i^{\mathrm c,\ast}(t),
\]
\[
\tilde u_{i,(k|t^+)}^{\mathrm c}-\tilde{\bar u}_i^{\mathrm c}(t^+)
=
u_{i,(k+1|t)}^{\mathrm c,\ast}-\bar u_i^{\mathrm c,\ast}(t).
\]
For the last stage, \eqref{eq:app_shift_x}, \eqref{eq:app_shift_u_terminal}, and \eqref{eq:app_terminal_eq} imply
\[
\tilde x_{i,(N_{\mathrm c}-1|t^+)}^{\mathrm c}
=
x_{i,(N_{\mathrm c}|t)}^{\mathrm c,\ast}
=
\bar x_i^{\mathrm c,\ast}(t),
\ 
\tilde u_{i,(N_{\mathrm c}-1|t^+)}^{\mathrm c}
=
\bar u_i^{\mathrm c,\ast}(t),
\]
so the last stage contributes
\(
\ell_i^{\mathrm c}(0,0)=0
\)
by Assumption~\ref{ass:regularity_cost}(3). Hence
\begin{equation}
\begin{aligned}
\tilde J_i^{\mathrm c}(t^+)
=
\sum_{k=1}^{N_{\mathrm c}-1}
&\ell_i^{\mathrm c}\!\Big(
x_{i,(k|t)}^{\mathrm c,\ast}-\bar x_i^{\mathrm c,\ast}(t),
u_{i,(k|t)}^{\mathrm c,\ast}-\bar u_i^{\mathrm c,\ast}(t)
\Big)
\\
&\;+\;
V_i^{\mathrm c}\!\bigl(\bar x_i^{\mathrm c,\ast}(t),x_i^{\mathrm{ref}}\bigr).
\end{aligned}
\label{eq:app_shifted_cost_identity_pre}
\end{equation}
Comparing \eqref{eq:app_shifted_cost_identity_pre} with the optimal contingency cost at time $t$,
\begin{equation}
\begin{aligned}
J_i^{\mathrm c,\ast}(t)
=
\sum_{k=0}^{N_{\mathrm c}-1}
&\ell_i^{\mathrm c}\!\Big(
x_{i,(k|t)}^{\mathrm c,\ast}-\bar x_i^{\mathrm c,\ast}(t),
u_{i,(k|t)}^{\mathrm c,\ast}-\bar u_i^{\mathrm c,\ast}(t)
\Big)
\\
&\;+\;
V_i^{\mathrm c}\!\bigl(\bar x_i^{\mathrm c,\ast}(t),x_i^{\mathrm{ref}}\bigr),
\end{aligned}
\label{eq:app_opt_cost_t}
\end{equation}
we obtain
\begin{equation}
\tilde J_i^{\mathrm c}(t^+)
=
J_i^{\mathrm c,\ast}(t)
-
\ell_i^{\mathrm c}\!\Big(
x_i(t)-\bar x_i^{\mathrm c,\ast}(t),
u_i(t)-\bar u_i^{\mathrm c,\ast}(t)
\Big),
\label{eq:app_shifted_cost_identity}
\end{equation}
because $x_i(t)=x_{i,(0|t)}^{\mathrm c,\ast}$ and $u_i(t)=u_{i,(0|t)}^{\mathrm c,\ast}$.
By \eqref{eq:lyap_bound_update_main}, the right-hand side of \eqref{eq:app_shifted_cost_identity} is exactly $\hat J_i^{\mathrm c}(t^+)$, which proves Lemma \ref{lem:appendix_shifted_cost_identity}.
\end{proof}

\section{Proof of Lemma~\ref{lem:disjoint_invariance}}
\begin{proof}
The claim is shown by induction.
Pairwise disjointness at $t=0$ holds by Assumption~\ref{ass:initial_feasibility}.
Assume that
\[
S_i^\ast(t)\cap S_j^\ast(t)=\emptyset,
\qquad \forall i\neq j,
\]
holds at some time $t\in\mathbb Z_+$.
Fix arbitrary $i\neq j$ and consider the update to $t^+$.

\smallskip
\noindent\textbf{Case 1: $\chi_i(t)=1$ and $\chi_j(t)=1$.}
Then
\(
S_i^\ast(t^+)=S_i^\ast(t),
\ 
S_j^\ast(t^+)=S_j^\ast(t),
\)
and hence
\[
S_i^\ast(t^+)\cap S_j^\ast(t^+)=\emptyset.
\]
\smallskip
\noindent\textbf{Case 2: $\chi_i(t)=1$ and $\chi_j(t)=0$.}
Then
\(
S_i^\ast(t^+)=S_i^\ast(t),
\
S_j^\ast(t^+)=\tilde S_j(t^+).
\)
Since $\chi_j(t)=0$, definition \eqref{eq:freezeIndicator} implies
\[
\tilde S_j(t^+)\cap S_i^\ast(t)=\emptyset.
\]
Therefore,
\[
S_i^\ast(t^+)\cap S_j^\ast(t^+)=\emptyset.
\]

\smallskip
\noindent\textbf{Case 3: $\chi_i(t)=0$ and $\chi_j(t)=1$.}
Symmetric to Case 2.

\smallskip
\noindent\textbf{Case 4: $\chi_i(t)=0$ and $\chi_j(t)=0$.}
Then
\(
S_i^\ast(t^+)=\tilde S_i(t^+),
\
S_j^\ast(t^+)=\tilde S_j(t^+).
\)
Since $\chi_i(t)=0$, definition \eqref{eq:freezeIndicator} implies
\[
\tilde S_i(t^+)\cap \tilde S_j(t^+)=\emptyset.
\]
Hence,
\[
S_i^\ast(t^+)\cap S_j^\ast(t^+)=\emptyset.
\]

Since the argument holds for every pair $i\neq j$, the family
$\{S_p^\ast(t^+)\}_{p\in\mathcal I}$ is pairwise disjoint.
This completes the induction.
\end{proof}

\section{Proof of Theorem~\ref{thm:rec_feas}}

\begin{proof}
We show that feasibility at time $t$ implies feasibility at time $t^+$.
Fix an arbitrary agent $i\in\mathcal I$ and suppose that the local FHOCP~\eqref{prob:local_fhocp} is feasible at time $t$.

By Lemma~\ref{lem:appendix_shifted_candidate_feasible}, the shifted contingency candidate defined by
\eqref{eq:app_shift_x}--\eqref{eq:app_shift_bar} satisfies all contingency-part constraints of \eqref{prob:local_fhocp} at time $t^+$.
By Lemma~\ref{lem:appendix_shifted_cost_identity}, its contingency cost satisfies
\[
\tilde J_i^{\mathrm c}(t^+)=\hat J_i^{\mathrm c}(t^+),
\]
so the Lyapunov-type constraint \eqref{prob:local_fhocp_lyap} is also feasible at time $t^+$.

Hence a feasible contingency candidate exists at time $t^+$.
By Assumption~\ref{ass:nominal_completion}, this feasible contingency candidate admits a feasible nominal completion satisfying the remaining constraints of the local FHOCP, including the shared-first-input constraint \eqref{prob:local_fhocp_shared}.
Therefore, the full local FHOCP~\eqref{prob:local_fhocp} is feasible at time $t^+$.

Since feasibility holds at $t=0$ by Assumption~\ref{ass:initial_feasibility}, induction yields feasibility for all $t\in\mathbb Z_+$.
\end{proof}

\section{Proof of Theorem~\ref{thm:safety}}

\begin{proof}
By Theorem~\ref{thm:rec_feas}, the local FHOCP remains feasible for every agent and every time.
Therefore the safe-set containment constraint \eqref{prob:local_fhocp_containment} holds along the closed loop, and thus
\(
\mathcal B_i(t)\subseteq S_i^\ast(t),
\ \forall i\in\mathcal I,\ \forall t\in\mathbb Z_+.
\)
By Lemma~\ref{lem:disjoint_invariance}, the active safe sets remain pairwise disjoint:
\[
S_i^\ast(t)\cap S_j^\ast(t)=\emptyset,
\qquad \forall i\neq j,\ \forall t\in\mathbb Z_+.
\]
Hence, for any distinct agents $i\neq j$,
\[
\mathcal B_i(t)\subseteq S_i^\ast(t),
\qquad
\mathcal B_j(t)\subseteq S_j^\ast(t),
\]
it holds that
\[
\mathcal B_i(t)\cap \mathcal B_j(t)=\emptyset \qquad \forall i\neq j,\ \forall t\in\mathbb Z_+.
\]
\end{proof}

\section{Proof of Proposition~\ref{prop:lyap_decrease}}

\begin{proof}
Fix an arbitrary agent $i\in\mathcal I$ and time $t\in\mathbb Z_+$.
By recursive feasibility, the local FHOCP~\eqref{prob:local_fhocp} is feasible at both $t$ and $t^+$.
By Lemma~\ref{lem:appendix_shifted_candidate_feasible}, the shifted contingency candidate is feasible at time $t^+$, and by Lemma~\ref{lem:appendix_shifted_cost_identity},
\[
\tilde J_i^{\mathrm c}(t^+)=\hat J_i^{\mathrm c}(t^+).
\]
Therefore, optimality of the solution at time $t^+$ gives
\begin{equation}
J_i^{\mathrm c,\ast}(t^+)
\le
\tilde J_i^{\mathrm c}(t^+)
=
\hat J_i^{\mathrm c}(t^+).
\label{eq:app_optimality_upper_bound}
\end{equation}
Substituting the bound update \eqref{eq:lyap_bound_update_main} into \eqref{eq:app_optimality_upper_bound} yields
\[
J_i^{\mathrm c,\ast}(t^+)
\le
J_i^{\mathrm c,\ast}(t)
-
\ell_i^{\mathrm c}\!\bigl(e_i(t),\delta u_i(t)\bigr),
\]
which is exactly \eqref{eq:lyap_decrease_main}.
Since $\ell_i^{\mathrm c}\ge 0$, the sequence $\{J_i^{\mathrm c,\ast}(t)\}_{t\in\mathbb Z_+}$ is monotonically nonincreasing.
Moreover, $J_i^{\mathrm c,\ast}(t)\ge 0$ for all $t$, because \eqref{eq:c_cost_def} is the sum of the nonnegative stage cost and the nonnegative offset cost.
Hence $\{J_i^{\mathrm c,\ast}(t)\}$ is lower bounded and therefore convergent.

Summing \eqref{eq:lyap_decrease_main} from $t=0$ to $T-1$ gives
\begin{equation}
\sum_{t=0}^{T-1}
\ell_i^{\mathrm c}\!\bigl(e_i(t),\delta u_i(t)\bigr)
\le
J_i^{\mathrm c,\ast}(0)-J_i^{\mathrm c,\ast}(T)
\le
J_i^{\mathrm c,\ast}(0).
\label{eq:app_telescoping_sum}
\end{equation}
Letting $T\to\infty$ yields \eqref{eq:summable_stage_main}.
In particular, the nonnegative sequence
$\ell_i^{\mathrm c}\!\bigl(e_i(t),\delta u_i(t)\bigr)$ converges to zero.

Finally, by Assumption~\ref{ass:regularity_cost}(3), there exists
$\underline{\alpha}_{\ell,i}\in\mathcal K_\infty$ such that
\[
\underline{\alpha}_{\ell,i}\!\bigl(\|(e_i(t),\delta u_i(t))\|\bigr)
\le
\ell_i^{\mathrm c}\!\bigl(e_i(t),\delta u_i(t)\bigr),
\qquad \forall t\in\mathbb Z_+.
\]
Suppose, for contradiction, that $\|(e_i(t),\delta u_i(t))\|\not\to 0$.
Then there exist $\varepsilon>0$ and an infinite subsequence $\{t_m\}_{m\in\mathbb Z_+}$ such that
\[
\|(e_i(t_m),\delta u_i(t_m))\|\ge \varepsilon,
\qquad \forall m\in\mathbb Z_+.
\]
By strict monotonicity of $\underline{\alpha}_{\ell,i}$,
\[
\ell_i^{\mathrm c}\!\bigl(e_i(t_m),\delta u_i(t_m)\bigr)
\ge
\underline{\alpha}_{\ell,i}(\varepsilon)
>0,
\qquad \forall m\in\mathbb Z_+,
\]
which contradicts
\[
\ell_i^{\mathrm c}\!\bigl(e_i(t),\delta u_i(t)\bigr)\to 0.
\]
Therefore,
\(
\|(e_i(t),\delta u_i(t))\|\to 0
\qquad\text{as } t\to\infty.
\)
\end{proof}

\section{Proof of Corollary~\ref{cor:fixed_ss_convergence}}

\begin{proof}
By Proposition~\ref{prop:lyap_decrease},
\[
\|(e_i(t),\delta u_i(t))\|\to 0.
\]
Assume that there exists $\bar t\in\mathbb Z_+$ such that
\[
\bar x_i^{\mathrm c,\ast}(t)=\bar x_i^{\mathrm c,\star},
\qquad
\bar u_i^{\mathrm c,\ast}(t)=\bar u_i^{\mathrm c,\star},
\qquad \forall t\ge \bar t.
\]
Then, $\forall t\ge \bar t$,
\(
e_i(t)=x_i(t)-\bar x_i^{\mathrm c,\star},
\
\delta u_i(t)=u_i(t)-\bar u_i^{\mathrm c,\star}.
\)
Hence
\[
x_i(t)\to \bar x_i^{\mathrm c,\star},
\qquad
u_i(t)\to \bar u_i^{\mathrm c,\star}
\qquad\text{as } t\to\infty.
\]
\end{proof}
\section{Proof of Corollary~\ref{cor:pnp_preservation}}

\begin{proof}
We consider join and leave events separately.

\smallskip
\noindent\textbf{Join.}
Let \(i_{\mathrm{new}}\notin\mathcal I\) request insertion at time
\(t_{\mathrm{join}}\). By Assumption~\ref{ass:pp_join_feasible}, the joining
agent can select an initial active safe set
\(S_{i_{\mathrm{new}}}^\ast(t_{\mathrm{join}})\) that is disjoint from all
representations \(\hat S_j(t_{\mathrm{join}})\) of the currently active
agents. If \(\chi_j(t_{\mathrm{join}})=0\), this representation equals the
true active safe set. If \(\chi_j(t_{\mathrm{join}})=1\),
Lemma~\ref{lem:pp_outer_approx_frozen} gives
\(S_j^\ast(t_{\mathrm{join}})\subseteq \hat S_j(t_{\mathrm{join}})\).
Hence,
\[
S_{i_{\mathrm{new}}}^\ast(t_{\mathrm{join}})
\cap
S_j^\ast(t_{\mathrm{join}})
=
\emptyset,
\qquad \forall j\in\mathcal I .
\]
Together with the feasibility of the joining agent and Theorem~\ref{thm:rec_feas}
for the previously active agents, the augmented system satisfies the
initial feasibility and separation conditions at \(t_{\mathrm{join}}\).
Applying Theorems~\ref{thm:rec_feas} and~\ref{thm:safety} from
\(t_{\mathrm{join}}\) onward proves recursive feasibility and collision
avoidance for the augmented system.

\smallskip
\noindent\textbf{Leave.}
If an agent leaves, it is removed from the locally reconstructed
constraints of the remaining agents. This only deletes constraints and
therefore cannot invalidate any previously feasible candidate solution.
The active safe sets of the remaining agents remain pairwise disjoint,
so recursive feasibility and collision avoidance follow again from
Theorems~\ref{thm:rec_feas} and~\ref{thm:safety}.
\end{proof}

\bibliographystyle{plain}
\bibliography{autosam}


\end{document}
